% !TeX program = xelatex
% 数学 LaTeX 排版 Skill v4.0.0；版式保持 v3 金标准. XeLaTeX 编译至少两次. 不分发字体.
% WZ-LOCKED-CLASS-BEGIN
\begin{filecontents*}[overwrite]{elegantbook-original-adapter.cls}
\NeedsTeXFormat{LaTeX2e}
\ProvidesClass{elegantbook-original-adapter}[2026/09/23 v3.0 fixed mathematical book environments]
\LoadClass[10pt,letterpaper,oneside]{book}
\RequirePackage[UTF8,scheme=plain,fontset=fandol]{ctex}
\RequirePackage{fontspec}
\setmainfont{cmunrm.otf}[BoldFont=cmunbx.otf,ItalicFont=cmunti.otf,BoldItalicFont=cmunbi.otf]
\setCJKfamilyfont{sourceheader}[ItalicFont=FandolKai-Regular.otf]{FandolKai-Regular.otf}
\RequirePackage{amsmath,amssymb,mathrsfs}
\RequirePackage{amsthm,etoolbox}
\RequirePackage{xcolor,geometry,setspace,fancyhdr,enumitem,needspace,xurl}
\definecolor{structurecolor}{HTML}{880E4F}
\definecolor{main}{HTML}{9C27B0}
\geometry{paperwidth=612bp,paperheight=792bp,left=20mm,right=20mm,top=95.04bp,bottom=20mm,headheight=12pt,headsep=25pt,footskip=30pt}
\setstretch{1.6}
\setlength{\parindent}{2em}
\setlength{\parskip}{0pt}
\setlist{nosep,leftmargin=2em}
\AtBeginDocument{\setlength{\abovedisplayskip}{4pt plus 1pt minus 1pt}\setlength{\belowdisplayskip}{4pt plus 1pt minus 1pt}\setlength{\abovedisplayshortskip}{2pt}\setlength{\belowdisplayshortskip}{3pt}}
\fancyhf{}
\fancyhead[L]{{\itshape\CJKfamily{sourceheader}\rightmark}}
\fancyhead[R]{{\itshape\CJKfamily{sourceheader}\leftmark}}
\fancyfoot[C]{\thepage}
\renewcommand{\headrulewidth}{0.4pt}
\renewcommand{\headrule}{{\color{structurecolor}\hrule height\headrulewidth width\headwidth\vskip-\headrulewidth}}
\pagestyle{fancy}
\fancypagestyle{cover}{\fancyhf{}\fancyfoot[C]{\thepage}\renewcommand{\headrulewidth}{0pt}\renewcommand{\headrule}{}}
\RequirePackage{hyperref}
\hypersetup{unicode=true,colorlinks=true,linkcolor=structurecolor,urlcolor=structurecolor,citecolor=structurecolor,linktoc=section,bookmarksnumbered=true,bookmarksopen=true,pdfborder={0 0 0},pdfstartview=FitH}
\setcounter{tocdepth}{1}
\renewcommand*\l@section{\@dottedtocline{1}{1.5em}{2.4em}}
\renewcommand{\thesection}{\thechapter.\arabic{section}}
\newcommand{\originalchapter}[1]{\par\addvspace{14pt}\Needspace{5\baselineskip}\refstepcounter{chapter}\setcounter{section}{0}\addcontentsline{toc}{chapter}{\protect\numberline{\thechapter}#1}\markboth{\thechapter\quad #1}{}\begingroup\centering\color{structurecolor}\bfseries\fontsize{14.4}{20}\selectfont\thechapter\quad #1\par\endgroup\nobreak\vspace{8pt}}
\newcommand{\originalsection}[1]{\par\addvspace{9pt}\Needspace{4\baselineskip}\refstepcounter{section}\addcontentsline{toc}{section}{\protect\hspace*{1.5em}\protect\numberline{\protect\textcolor{structurecolor}{\thesection}}#1}\markright{\thesection\quad #1}\noindent{\fontsize{12}{17}\selectfont\color{structurecolor}\thesection\quad}{\bfseries\fontsize{12}{17}\selectfont #1}\par\nobreak\vspace{4pt}}

% Semantic environments are registered once here, not re-designed in manuscripts.
\newtheoremstyle{wzstatement}{6pt}{5pt}
  {\normalfont\kaishu}{0pt}{\normalfont\bfseries\color{structurecolor}}
  {}{.7em}{\thmname{#1}\thmnumber{ #2}\thmnote{\normalfont\ (#3)}}
\newtheoremstyle{wzdefinition}{6pt}{5pt}
  {\normalfont\songti}{0pt}{\normalfont\bfseries\color{structurecolor}}
  {}{.7em}{\thmname{#1}\thmnumber{ #2}\thmnote{\normalfont\ (#3)}}
\newtheoremstyle{wzproblem}{7pt}{5pt}
  {\normalfont\songti}{2em}{\normalfont\bfseries\color{main}}
  {}{.7em}{\thmname{#1}\thmnumber{ #2}}
\newtheoremstyle{wzremark}{4pt}{4pt}
  {\normalfont\songti}{0pt}{\normalfont\bfseries}
  {}{.7em}{\thmname{#1}}
\theoremstyle{wzstatement}
\newtheorem{theorem}{定理}[chapter]
\newtheorem{lemma}[theorem]{引理}
\newtheorem{proposition}[theorem]{命题}
\newtheorem{corollary}[theorem]{推论}
\theoremstyle{wzdefinition}
\newtheorem{definition}[theorem]{定义}
\theoremstyle{wzproblem}
\newtheorem{example}{例题}[chapter]
\newtheorem{exercise}{习题}[chapter]
\theoremstyle{wzremark}
\newtheorem*{remark}{注}
\renewcommand{\proofname}{证明}
% Retain the native amsthm QED stack; only adjust the fixed typography.
\expandafter\patchcmd\csname\string\proof\endcsname{\normalfont \topsep6\p@\@plus6\p@\relax}
  {\normalfont\songti \topsep5\p@\@plus1\p@\relax}
  {}{\ClassError{elegantbook-original-adapter}{Cannot patch proof spacing}{Check the amsthm version.}}
\expandafter\patchcmd\csname\string\proof\endcsname{\itshape}{\normalfont\bfseries}
  {}{\ClassError{elegantbook-original-adapter}{Cannot patch proof heading}{Check the amsthm version.}}
\expandafter\patchcmd\csname\string\proof\endcsname{\ignorespaces}
  {\setlength{\parindent}{2em}\ignorespaces}
  {}{\ClassError{elegantbook-original-adapter}{Cannot patch proof paragraphs}{Check the amsthm version.}}
\AtBeginEnvironment{proof}{\par\Needspace{3\baselineskip}}
\renewcommand{\qedsymbol}{\ensuremath{\square}}
\newenvironment{solution}{\begin{proof}[解答]}{\end{proof}}
\newcommand{\wzcheckchapter}{\ifnum\value{chapter}<1
  \ClassError{elegantbook-original-adapter}{Numbered mathematics before chapter 1}{Move it after the first originalchapter.}\fi}

\newcommand{\wzregisterguard}[1]{\AtBeginEnvironment{#1}{\wzcheckchapter\par\Needspace{3\baselineskip}}}
\forcsvlist{\wzregisterguard}{theorem,lemma,proposition,corollary,definition,example,exercise}
% Front matter and bibliography are not numbered mathematical chapters.
\newcommand{\originalpreface}{\par\addvspace{10pt}\Needspace{5\baselineskip}
  \noindent{\bfseries\fontsize{12}{17}\selectfont 前言}\par\nobreak\vspace{4pt}}
\newcommand{\originalcontents}{\par\addvspace{14pt}
  \begin{center}{\color{structurecolor}\bfseries\fontsize{14.4}{20}\selectfont 目录}\end{center}
  \vspace{-10pt}\@starttoc{toc}}
\renewcommand{\bibname}{参考文献}
\renewenvironment{thebibliography}[1]{
  \par\addvspace{12pt}\Needspace{3\baselineskip}\phantomsection
  \addcontentsline{toc}{chapter}{\bibname}
  \markboth{\bibname}{}
  \noindent{\color{structurecolor}\bfseries\fontsize{14.4}{20}\selectfont\bibname}\par\nobreak\vspace{6pt}
  \list{\@biblabel{\@arabic\c@enumiv}}
    {\settowidth\labelwidth{\@biblabel{#1}}\leftmargin\labelwidth
     \advance\leftmargin\labelsep\usecounter{enumiv}
     \let\p@enumiv\@empty\renewcommand\theenumiv{\@arabic\c@enumiv}
     \setlength{\itemsep}{3pt}\setlength{\parsep}{0pt}}
  \sloppy\clubpenalty4000\widowpenalty4000\sfcode`\.=1000\relax
}{\def\@noitemerr{\@latex@warning{Empty bibliography}}\endlist}

\numberwithin{equation}{chapter}
\renewcommand{\theequation}{\thechapter.\arabic{equation}}
\allowdisplaybreaks[2]
\emergencystretch=2em
\clubpenalty=6000
\widowpenalty=6000
\raggedbottom
\endinput
\end{filecontents*}
% WZ-LOCKED-CLASS-END
\documentclass{elegantbook-original-adapter}
% ===== 元数据内容槽 =====
\newcommand{\BookTitle}{复变函数与应用}
\newcommand{\BookAuthor}{Jeremy Orloff 原讲义; 中文译编与证明补充}
\newcommand{\BookDate}{MIT 18.04, Spring 2018; 18.112, Fall 2008 拓展}
\hypersetup{pdftitle={复变函数与应用: MIT OCW 中文重编本},pdfauthor={Jeremy Orloff; Sigurdur Helgason; Zuoqin Wang; 中文译编}}
\usepackage{tikz,pgfplots,booktabs,longtable}
\pgfplotsset{compat=1.18}
\usetikzlibrary{arrows.meta,calc,positioning,decorations.markings}
\newcommand{\C}{\mathbb C}
\newcommand{\R}{\mathbb R}
\newcommand{\N}{\mathbb N}
\newcommand{\Z}{\mathbb Z}
\newcommand{\D}{\mathbb D}
\newcommand{\dd}{\,\mathrm d}
\newcommand{\ii}{\mathrm i}
\newcommand{\ee}{\mathrm e}
\newcommand{\eps}{\varepsilon}
\newcommand{\Res}{\operatorname{Res}}
\newcommand{\Ind}{\operatorname{Ind}}
\newcommand{\Log}{\operatorname{Log}}
\newcommand{\Arg}{\operatorname{Arg}}
\newcommand{\dist}{\operatorname{dist}}
\newcommand{\sourceinfo}[1]{\par\noindent\textit{来源: #1.}\par}
\begin{document}
\frontmatter
\begin{center}
{\color{structurecolor}\bfseries\fontsize{18}{25}\selectfont\BookTitle\par}
{\bfseries MIT OCW 中文重编本\par}
\BookAuthor\par
\BookDate\par
\end{center}
\originalpreface
复可微要求从平面内所有方向趋近时得到同一个导数. 这一条件使局部微分、全局积分和幂级数表达相互联系. 本书从复数与初等函数开始, 建立 Cauchy 理论、级数与留数方法, 再进入辐角原理、调和函数、保角映射、流体势、Laplace 变换与解析延拓. 第14章讨论正规族与 Riemann 映射定理, 属于拓展.

主源为 Jeremy Orloff 的 \href{https://ocw.mit.edu/courses/18-04-complex-variables-with-applications-spring-2018/pages/lecture-notes/}{MIT 18.04 Complex Variables with Applications, Spring 2018}: 实际15份文件, 包括13个专题、课程介绍和多元微积分复习, 共249个原件物理页. 课程介绍注明课程设计源于 Andre Nachbin, Jörn Dunkel 参与校订. 拓展依据 \href{https://ocw.mit.edu/courses/18-112-functions-of-a-complex-variable-fall-2008/pages/lecture-notes/}{MIT 18.112 Functions of a Complex Variable, Fall 2008}, 由 Sigurdur Helgason 授课, Zuoqin Wang 在其指导下整理补充讲义. 实际22份补充文件共76页, 并非一本完整的 Ahlfors 教材. 本书不归档或翻译课程指定的第三方教材.

读者先修为一元微积分、实幂级数、二元微分与积分的基本知识. 实数完备性、紧集上连续函数的极值与一致连续性、实微积分基本定理和有界区域的 Green 定理作为先修工具. 使用一致收敛、局部紧性与 Fourier 反演时, 正文会补出所需版本及其依据. 主源中的图像全部用原生 TikZ 重绘或用文字推导替代; 独立习题集、习题课和考试没有整批纳入, 每章配套题及解答由编者依本章目标另设.

这是按数学依赖关系译编的完整主课正文, 不逐句保留英文语序, 不逐页复刻排版. 重复定义、复习段及官方署名尾页在中文正文中合并, 原始文件保持不动. 为消除原课略去的技术细节, Goursat、同伦不变性、Laurent 展开、局部映射、积分反演等论证均有编者补充, 在相应位置明示. Gamma 的 Stirling 渐近采用实正轴版本, 不把它误写为无条件的全复平面近似.

原讲义与中文改编按 \href{https://creativecommons.org/licenses/by-nc-sa/4.0/}{CC BY-NC-SA 4.0 International} 非商业使用及相同方式共享. 署名对象为 Massachusetts Institute of Technology、原课程与上述作者. 翻译、章节重组、重绘、证明补足和自设练习属于本版改动, 不表示 MIT 或原作者认可本项目. 单独权利声明优先于默认许可. 官方条款见 \href{https://ocw.mit.edu/pages/privacy-and-terms-of-use/}{MIT OCW 使用条款}. 逐文件下载地址、页数、哈希、许可核查和覆盖范围见书末来源表及源码中的 source-manifest.json. 核查日期为2026-10-08.
\clearpage
\originalcontents
\clearpage
\mainmatter
\originalchapter{复数与复平面}
\sourceinfo{18.04 Topic 1, pp.1--24; Topic 2, pp.5--7. 球面坐标推导为编者补充}
\originalsection{代数与几何}
\begin{definition}
复数是形如 $z=x+\ii y$ 的数, 其中 $x,y\in\R$, $\ii^2=-1$. 记 $\Re z=x$, $\Im z=y$, $\bar z=x-\ii y$, $|z|=(x^2+y^2)^{1/2}$. 全体复数组成域 $\C$.
\end{definition}
加法按两个坐标相加, 乘法按分配律与 $\ii^2=-1$ 计算. $\Im z$ 是实数, 不是 $\ii y$. 恒等式 $z\bar z=|z|^2$ 给出非零数的逆元 $z^{-1}=\bar z/|z|^2$. 例如 $(3+4\ii)/(1+2\ii)=(11-2\ii)/5$. 共轭满足 $\overline{z+w}=\bar z+\bar w$, $\overline{zw}=\bar z\bar w$; 因而 $|zw|=|z||w|$.

把 $z$ 对应到点 $(x,y)$, 模就是到原点的距离, $|z-w|$ 就是两点距离. 平移由加法表示, 复数乘法的旋转意义将在极坐标中出现.
\begin{proposition}[三角不等式]
任意 $z,w\in\C$ 满足 $|z+w|\le |z|+|w|$ 和 $\bigl||z|-|w|\bigr|\le |z-w|$. 第一式取等当且仅当至少一个数为零, 或两个非零数在原点发出的同一条射线上.
\end{proposition}
\begin{proof}
写 $z\bar w=A+\ii B$, 则 $A\le(A^2+B^2)^{1/2}=|z||w|$. 于是 $|z+w|^2=|z|^2+2A+|w|^2\le(|z|+|w|)^2$. 两边非负, 开方即可. 非零时取等要求 $z\bar w$ 为正实数, 即 $z/w$ 为正实数. 再对 $z=(z-w)+w$ 应用第一式, 得 $|z|-|w|\le|z-w|$; 交换 $z,w$ 得另一半.
\end{proof}
圆和直线可以直接写成复数关系. 圆 $|z-a|=r$ 的内部为 $|z-a|<r$; 直线 $\Re(\bar b z)=c$ 的法向量为 $b\ne0$. 轨迹 $|z-a|=|z-b|$ 是线段 $ab$ 的垂直平分线. 这些表达式不需要把复数大小作有序比较; $z>w$ 在复数域内没有通常的意义.

\originalsection{指数与极坐标}
对实数 $\theta$, 定义 $\ee^{\ii\theta}=\cos\theta+\ii\sin\theta$. 实三角函数的加法公式给出 $\ee^{\ii\theta}\ee^{\ii\phi}=\ee^{\ii(\theta+\phi)}$. 对 $z=x+\ii y$ 定义 $\ee^z=\ee^x(\cos y+\ii\sin y)$, 于是 $\ee^{z+w}=\ee^z\ee^w$, $\ee^0=1$, $|\ee^z|=\ee^{\Re z}$, 指数永远不为零. 实指数、正弦与余弦的级数绝对收敛, 分离偶项和奇项可得 $\ee^z=\sum_{n\ge0}z^n/n!$. 第5章会证明这一复幂级数可以逐项求导.

每个非零复数可写成 $z=r\ee^{\ii\theta}$, 其中 $r=|z|>0$. 所有满足该式的 $\theta$ 构成辐角集合 $\arg z=\theta+2\pi\Z$. 两个数相乘时模相乘、辐角相加; 除法则模相除、辐角相减. 因而乘以 $2\ii$ 将所有长度放大两倍, 同时逆时针转过 $\pi/2$.
\begin{center}
\begin{tikzpicture}[scale=1.0,>=Stealth]
\draw[->] (-.4,0)--(4,0) node[right]{$\Re z$};
\draw[->] (0,-.3)--(0,2.8) node[above]{$\Im z$};
\draw[thick,->] (0,0)--(3,2) node[above right]{$z=r\ee^{\ii\theta}$};
\draw[dashed] (3,0) node[below]{$x$}--(3,2)--(0,2) node[left]{$y$};
\draw (0.8,0) arc (0:33.69:.8); \node at (1,.28) {$\theta$};
\node at (1.5,1.3) {$r$};
\end{tikzpicture}
\end{center}
\begin{example}
求 $(1+\ii)^6$ 及 $w^n=z\ne0$ 的全部解, 其中 $n$ 为正整数.
\end{example}
\begin{solution}
$1+\ii=\sqrt2\ee^{\ii\pi/4}$, 因而 $(1+\ii)^6=8\ee^{3\pi\ii/2}=-8\ii$. 若 $z=r\ee^{\ii\theta}$, 所有根为 $r^{1/n}\ee^{\ii(\theta+2k\pi)/n}$, $k=0,\ldots,n-1$. 任意根的模必须为 $r^{1/n}$, 辐角必须满足 $n\phi-\theta\in2\pi\Z$, 所以已经找全; 不同 $k$ 给出不同根. 它们是圆上正 $n$ 边形的顶点.
\end{solution}
根的多值性不能通过把辐角默认为一个数而消失. 取某个连续辐角范围, 就是在选根函数的分支. 例如正实轴上取正平方根, 沿绕原点一周的曲线连续延伸后会变成负平方根.

复三角函数定义为 $\sin z=(\ee^{\ii z}-\ee^{-\ii z})/(2\ii)$, $\cos z=(\ee^{\ii z}+\ee^{-\ii z})/2$. 直接展开得到
\[
\sin(x+\ii y)=\sin x\cosh y+\ii\cos x\sinh y,\qquad
\cos(x+\ii y)=\cos x\cosh y-\ii\sin x\sinh y.
\]
加法公式和 $\sin^2z+\cos^2z=1$ 在复数上仍成立, 但 $|\sin z|\le1$ 不成立: $\sin(\ii y)=\ii\sinh y$. 解 $\sin z=0$ 等价于 $\ee^{2\ii z}=1$, 由模先得 $\Im z=0$, 再得 $z\in\pi\Z$. $\cos z=0$ 的解为 $\pi/2+\pi\Z$.

\originalsection{对数与映射}
指数的周期为 $2\pi\ii$, 且 $\ee^{z_1}=\ee^{z_2}$ 当且仅当 $z_1-z_2\in2\pi\ii\Z$. 因此方程 $\ee^w=z\ne0$ 的解为 $w=\log|z|+\ii(\theta+2k\pi)$. 这里 $\log|z|$ 是实对数. 在割平面 $\C\setminus(-\infty,0]$ 上选择 $-\pi<\Arg z<\pi$, 定义主值 $\Log z=\log|z|+\ii\Arg z$. 对负实数可给出数值主辐角 $\pi$, 但这样得到的函数在该轴附近不连续, 不能当作该轴邻域的全纯分支.

在选定对数分支的域上可定义 $z^\alpha=\ee^{\alpha\Log z}$. 即使 $\alpha$ 为实数, 跨越割线后数值也可能改变. 主值通常不满足 $\Log(zw)=\Log z+\Log w$: 例如 $z=w=\ee^{3\pi\ii/4}$, 两边分别为 $-\pi\ii/2$ 与 $3\pi\ii/2$.

复函数也描述平面区域的变换. $w=z^2$ 把辐角为 $(0,\pi/2)$ 的第一象限一一映为上半平面; 模平方而辐角加倍. $w=\ee^z$ 把带状域 $0<\Im z<\pi$ 一一映为上半平面: 每条水平线变成射线, 每条竖直线段变成半圆弧. 逆映射是辐角在 $(0,\pi)$ 的对数分支. 区域的一一对应须同时核对像、单射和满射, 只检查边界图形并不够.

\originalsection{无穷远点}
扩充复平面为 $\widehat\C=\C\cup\{\infty\}$. 定义 $z_n\to\infty$ 为 $|z_n|\to\infty$, 无论趋近方向如何都到同一个点. 无穷远邻域为 $\{\infty\}\cup\{|z|>R\}$. 局部坐标 $\zeta=1/z$ 将它变为零点邻域; 研究 $f$ 在无穷远的极限、全纯性或奇点, 就研究 $f(1/\zeta)$ 在零点的相应性质.

单位球面上的点 $(X,Y,T)$ 经北极投影到 $z=(X+\ii Y)/(1-T)$. 反解为
\[
X=\frac{2\Re z}{1+|z|^2},\qquad Y=\frac{2\Im z}{1+|z|^2},\qquad
T=\frac{|z|^2-1}{1+|z|^2}.
\]
把这些式子代入 $X^2+Y^2+T^2=1$ 可核实球面条件. 当 $|z|\to\infty$, 球面点趋近北极 $(0,0,1)$, 这解释了只添加一个无穷远点的几何意义. 非常数多项式的无穷远极限是无穷大, 而 $\ee^z$ 在无穷远没有单一极限.

\originalsection{习题与解答}
\begin{exercise}
求 $z^4=-16$ 的全部根, 并证明四根之和为零.
\end{exercise}
\begin{solution}
根为 $2\ee^{\ii(\pi/4+k\pi/2)}$, $k=0,1,2,3$. 相隔两个指标的根互为相反数, 所以和为零. 也可用 $1+\ii-1-\ii=0$ 求和.
\end{solution}
\begin{exercise}
证明 $|\sin(x+\ii y)|^2=\sin^2x+\sinh^2y$, 并求 $0\le x,y\le2\pi$ 上 $|\sin z|$ 的最大值及取等位置.
\end{exercise}
\begin{solution}
利用实虚部表达及 $\cosh^2y-\sinh^2y=1$, 有 $\sin^2x\cosh^2y+\cos^2x\sinh^2y=\sin^2x+\sinh^2y$. 两项分别在 $x=\pi/2,3\pi/2$ 和 $y=2\pi$ 取最大值, 模的最大值为 $\cosh(2\pi)$, 位置即这两个边界点.
\end{solution}
\begin{exercise}
计算 $\int\ee^x\cos2x\dd x$.
\end{exercise}
\begin{solution}
取 $\int\ee^{(1+2\ii)x}\dd x=\ee^{(1+2\ii)x}/(1+2\ii)$ 的实部. 乘分子分母以 $1-2\ii$, 得 $\ee^x(\cos2x+2\sin2x)/5+C$. 取实部与实积分交换是积分线性的直接结果.
\end{solution}

\originalchapter{全纯函数}
\sourceinfo{18.04 Topic 2, pp.1--20. 偏导条件与单点可微的区别为编者补充}
\originalsection{极限与复导数}
\begin{definition}
记 $D(a,r)=\{z:|z-a|<r\}$. 若集合中每点都有包含于集合的小圆盘, 则它为开集. 连通开集称为域. 极限 $f(z)\to L$ ($z\to a$) 指: 对每个 $\eps>0$, 存在 $\delta>0$, 使定义域中 $0<|z-a|<\delta$ 时 $|f(z)-L|<\eps$. 连续指此极限等于 $f(a)$.
\end{definition}
该条件涵盖所有趋近方式. 验证有限条直线上的极限相同, 不能证明二维极限存在. 实部和虚部极限分别存在且等于 $\Re L,\Im L$, 才等价于复极限存在. 这是因为 $|\Re w|,|\Im w|\le |w|\le |\Re w|+|\Im w|$.
\begin{definition}
若有限极限 $f'(a)=\lim_{h\to0}(f(a+h)-f(a))/h$ 存在, 称 $f$ 在 $a$ 复可微. 在开集内每点复可微称为全纯. 局部有收敛幂级数表达称为解析; 第5章证明这两个性质等价. 在整个 $\C$ 全纯称为整函数.
\end{definition}
原课程常把单点复可微也称作 ``analytic at a point''. 本书保留上述区分, 防止把一个点的导数误当作邻域内的全纯性. 导数存在等价于 $f(a+h)=f(a)+f'(a)h+o(|h|)$, 因而复可微蕴含连续.

$z^2$ 的差商为 $2z+h$, 所以导数为 $2z$. 共轭函数的差商为 $\bar h/h$: 实方向得1, 虚方向得 $-1$, 因而处处没有复导数. 函数 $|z|^2$ 在0的差商为 $\bar h\to0$, 因而在0可微, 却在任何零点邻域内都不全纯.

和、积、商及链式法则与实微积分形式相同. 其共同依据是一次近似. 例如乘积展开为
\[
(fg)(a+h)=f(a)g(a)+\bigl(f'(a)g(a)+f(a)g'(a)\bigr)h+o(|h|).
\]
商法则要求分母在点处非零. 链式法则中先把内函数的增量代入外函数的一次近似, 利用内增量为 $O(|h|)$ 控制余项, 得 $(g\circ f)'=g'(f)f'$. 不能因为熟悉实公式而忽略这些定义域条件.

\originalsection{Cauchy--Riemann 方程}
\begin{theorem}\label{thm:cr}
写 $f=u+\ii v$, $z=x+\ii y$. 若 $f$ 在点 $a$ 复可微, 则该点四个偏导存在且 $u_x=v_y$, $u_y=-v_x$, $f'=u_x+\ii v_x$. 反之, 若 $u,v$ 在该点作为实二元函数可微且满足两条方程, 则 $f$ 在该点复可微. 特别地, 邻域内四个偏导连续并满足方程, 则 $f$ 在该邻域全纯.
\end{theorem}
\begin{proof}
差商沿实方向趋近时为 $u_x+\ii v_x$, 沿虚方向为 $v_y-\ii u_y$. 两者相等给出必要条件. 反向论证须控制任意方向. 实可微给出
\[
\Delta f=(u_x+\ii v_x)\Delta x+(u_y+\ii v_y)\Delta y+o\bigl((\Delta x^2+\Delta y^2)^{1/2}\bigr).
\]
代入方程, 右端为 $(u_x+\ii v_x)(\Delta x+\ii\Delta y)+o(|\Delta z|)$. 除以 $\Delta z$ 后余项趋零, 得到复导数. 连续偏导蕴含实可微是二元微积分的结论.
\end{proof}
仅在一点有偏导并满足两条方程还不充分. 设 $f(0)=0$, 其余点 $f(x+\ii y)=xy/(x^2+y^2)$. 在零点四个偏导全为零, 但沿 $y=x$ 得 $f=1/2$, 函数甚至不连续. 判定全纯时, 必须检查正则性, 再检查方程.

实导数矩阵在全纯点形如 $\begin{pmatrix}a&-b\\b&a\end{pmatrix}$, 其中 $a+\ii b=f'(z)$. 它等于一次旋转乘上一个伸缩; 行列式为 $|f'(z)|^2$. 当 $f'\ne0$ 时, 局部一阶映射保持有向角. 第11章会把这种一阶结论发展成局部和全局的保角映射.
\begin{proposition}
域上全纯函数若导数恒为零, 则为常数. 若其值全为实数, 或实部恒定, 或模恒定, 也必为常数.
\end{proposition}
\begin{proof}
导数为零时由方程知 $u,v$ 四个偏导为零. 在任意小圆盘上沿线段应用实微积分, 两者恒定. 局部常数在连通域上恒定: 固定一个函数值, 其对应点集及补集都为开集. 纯实值与实部恒定的情形直接由方程得到四个偏导为零. 模恒定为 $c>0$ 时, 微分 $u^2+v^2=c^2$ 并代入方程, 得关于 $u_x,v_x$ 的线性方程, 系数行列式为 $u^2+v^2=c^2$, 因而导数为零. $c=0$ 时已经恒为零.
\end{proof}

\originalsection{初等函数与分支}
由 $\ee^{x+\ii y}=\ee^x\cos y+\ii\ee^x\sin y$ 检查方程, 得 $(\ee^z)'=\ee^z$. 三角函数由指数线性组合得到 $(\sin z)'=\cos z$, $(\cos z)'=-\sin z$. 多项式整, 有理函数在分母非零处全纯, 并不能把未约掉的零分母全部称为极点.

主对数在割平面全纯. 局部极坐标微分给出 $(\log r)_x=x/r^2$, $(\log r)_y=y/r^2$, $\theta_x=-y/r^2$, $\theta_y=x/r^2$, 符合方程, 所以 $(\Log z)'=(x-\ii y)/r^2=1/z$. 任意两条全纯对数分支在连通共同域上的差为 $2k\pi\ii$, 其中整数 $k$ 恒定. 这是连续函数只能把连通集映为连通集的结果.

没有在 $\C\setminus\{0\}$ 上定义的连续全局对数. 若有, 对单位圆参数 $\ee^{\ii t}$ 可写 $L(\ee^{\ii t})=\ii t+2\pi\ii k(t)$. 连续性使整数函数 $k$ 恒定, 但 $t=0,2\pi$ 是同一点却给出相差 $2\pi\ii$ 的值, 矛盾.
\begin{example}
在割平面上求 $z^\alpha$ 的导数, 并解释 $\sqrt z$ 在0处的行为.
\end{example}
\begin{solution}
固定 $L=\Log$, 链式法则给出 $(\ee^{\alpha L})'=\alpha\ee^{\alpha L}/z=\alpha z^{\alpha-1}$, 右边须使用同一分支. 平方根分支在割平面全纯, 可连续补为 $\sqrt0=0$, 但零点差商的模为 $1/\sqrt{|h|}$, 不会有有限极限. 零点是分支点, 不是后文意义下的孤立全纯奇点.
\end{solution}

\originalsection{双曲函数与复合割线}
双曲函数定义为 $\cosh z=(\ee^z+\ee^{-z})/2$, $\sinh z=(\ee^z-\ee^{-z})/2$. 它们整, 导数互换, 且 $\cosh^2z-\sinh^2z=1$, $\cos(\ii z)=\cosh z$, $\sin(\ii z)=\ii\sinh z$. 这些恒等式从指数代数直接导出, 因而不受实数自变量的限制.

反三角函数必须处理多个分支. 解 $w=\sin z$, 令 $\zeta=\ee^{\ii z}\ne0$, 得 $\zeta^2-2\ii w\zeta-1=0$, 所以 $\zeta=\ii w\pm\sqrt{1-w^2}$, 再取 $z=-\ii\log\zeta$. 在 $w=0$ 附近选择平方根在0为1和对数在1为0, 得
\[
\arcsin w=-\ii\Log\bigl(\ii w+\sqrt{1-w^2}\bigr).
\]
该式是局部解析逆, 不是未经选择的全局单值公式. 隐式微分给 $(\arcsin w)'=1/\sqrt{1-w^2}$, 根的符号要与初值相容. $w=\pm1$ 是分支点. 类似地, 解 $w=\sinh z$ 令 $\zeta=\ee^z$, 得局部 $\operatorname{arsinh}w=\Log(w+\sqrt{w^2+1})$, 在0取导数1的分支. 多值对数与平方根的其他选择给其余逆值.

复合函数的割线是外函数割线的原像. 主平方根 $\sqrt{1-z}$ 的域为 $\C\setminus[1,\infty)$: 条件 $1-z\in(-\infty,0]$ 恰等价于 $z$ 在该射线上. 对主值 $\sqrt{1+\ee^z}$, 写 $z=x+\ii y$. 外根被割掉当且仅当 $1+\ee^z$ 为非正实数. 这要求 $\sin y=0$; 偶数倍 $\pi$ 时值为正, 奇数倍时值 $1-\ee^x\le0$ 当且仅当 $x\ge0$. 所以须删去全部水平射线 $\{x+(2k+1)\pi\ii:x\ge0\}$, $k\in\Z$. 在删去这些射线后的开集, 指数与主平方根的复合全纯. 只删某一个零点, 或只删负实轴, 都没有检查复合割线的实际原像.

\originalsection{习题与解答}
\begin{exercise}
找出 $f(z)=|z|^2$ 的全部复可微点, 以及 $f(z)=\Re z$ 的全部复可微点.
\end{exercise}
\begin{solution}
前者 $u=x^2+y^2$, $v=0$, 方程要求 $x=y=0$, 且已由差商验证0处导数为0. 后者 $u_x=1$, $v_y=0$, 必要条件处处失败, 所以无点复可微.
\end{solution}
\begin{exercise}
设 $u=x^3-3xy^2+2x$. 求 $v$ 使 $u+\ii v$ 全纯.
\end{exercise}
\begin{solution}
$v_y=u_x=3x^2-3y^2+2$, 积分得 $v=3x^2y-y^3+2y+c(x)$. 再由 $v_x=-u_y=6xy$ 得 $c'(x)=0$. 故 $v=3x^2y-y^3+2y+C$, 函数为 $z^3+2z+\ii C$.
\end{solution}
\begin{exercise}
设全纯 $f=u+\ii v$, 在 $r>0$ 的局部极坐标中推导 Cauchy--Riemann 方程.
\end{exercise}
\begin{solution}
链式法则给 $u_r=u_x\cos\theta+u_y\sin\theta$ 和 $v_\theta=r(-v_x\sin\theta+v_y\cos\theta)$. 代入直角坐标方程得 $u_r=v_\theta/r$. 同理 $v_r=-u_\theta/r$. 逆向用坐标变换可恢复原方程; 在 $r=0$ 不能使用这组表达.
\end{solution}

\originalchapter{复积分与 Cauchy 定理}
\sourceinfo{18.04 Topic 3 全部正文; 多元微积分复习. Goursat 与同伦证明为编者补充, 参照18.112 L8--L9的范围}
\originalsection{路径积分}
\begin{definition}
路径 $\gamma:[a,b]\to\C$ 是连续、分段 $C^1$ 的映射, 允许自交和重复经过同一点. 对其像上连续的函数 $f$, 定义 $\int_\gamma f(z)\dd z=\int_a^b f(\gamma(t))\gamma'(t)\dd t$. 长度为 $\ell(\gamma)=\int_a^b|\gamma'(t)|\dd t$.
\end{definition}
保向的分段光滑重参数化不改变积分, 由实换元公式得证. 反向路径记为 $-\gamma$, 积分改变符号; 首尾可衔接的路径拼接后积分相加. $\dd z$ 表示有向复增量, $|\dd z|$ 表示无向弧长. 若 $f=u+\ii v$, 则
\[
\int_\gamma f\dd z=\int_\gamma(u\dd x-v\dd y)+\ii\int_\gamma(v\dd x+u\dd y).
\]
这与两个实线积分完全一致, 却把两组微分同时组织起来.
\begin{proposition}[积分估计]\label{prop:ml}
有 $|\int_\gamma f\dd z|\le\int_\gamma|f||\dd z|\le M\ell(\gamma)$, 其中 $M=\max_\gamma|f|$.
\end{proposition}
\begin{proof}
复数 Riemann 和满足三角不等式. 取极限得 $|\int_a^b g(t)\dd t|\le\int_a^b|g(t)|\dd t$. 代入 $g=f(\gamma)\gamma'$ 得第一式, 再用 $|f|\le M$ 得第二式.
\end{proof}
\begin{example}
在逆时针圆 $z=a+r\ee^{\ii t}$ 上计算 $\int(z-a)^n\dd z$, 其中 $n\in\Z$.
\end{example}
\begin{solution}
代入参数得到 $\ii r^{n+1}\int_0^{2\pi}\ee^{\ii(n+1)t}\dd t$. $n=-1$ 时为 $2\pi\ii$, 其余整数时为0. 这一结果是 Cauchy 公式和留数理论的基本计算.
\end{solution}
同一个单位圆上 $\int\bar z\dd z=\int_0^{2\pi}\ee^{-\ii t}\ii\ee^{\ii t}\dd t=2\pi\ii$. 因而闭路积分为零并非所有连续函数的性质.

\originalsection{原函数与路径无关}
\begin{theorem}
设 $f$ 在域 $\Omega$ 连续. 下列三项等价: $f$ 在 $\Omega$ 有全纯原函数 $F$, 即 $F'=f$; 同端点路径积分相同; 所有闭路径积分为零.
\end{theorem}
\begin{proof}
若 $F'=f$, 链式法则与实微积分给 $\int_\gamma f\dd z=F(\gamma(b))-F(\gamma(a))$. 因而原函数蕴含路径无关, 并蕴含闭路积分为零. 反过来, 两条同端点路径与反向路径拼接成闭路, 所以闭路为零蕴含路径无关. 连通开集内任意两点可用折线连接: 固定基点, 可连到它的点集及其补集都开. 定义 $F(z)=\int_{z_0}^z f\dd\zeta$, 因路径无关而良定义. 对足够小的复数 $h$, 选择末段为直线, 得
\[
\frac{F(z+h)-F(z)}h=\int_0^1 f(z+th)\dd t\longrightarrow f(z).
\]
收敛由 $f$ 在 $z$ 的连续性一致控制 $0\le t\le1$ 而得到. 所以 $F'=f$.
\end{proof}
例如 $1/z^2$ 在穿孔平面有原函数 $-1/z$, 所有闭路积分为零; $1/z$ 没有全局原函数, 因为单位圆积分不为零. 函数全纯与域没有洞, 是两个不同条件.

\originalsection{三角形积分定理}
\begin{theorem}[Goursat]\label{thm:goursat}
若 $f$ 在闭三角形 $T$ 的某开邻域全纯, 则 $\int_{\partial T}f\dd z=0$. 不需要预先假定 $f'$ 连续.
\end{theorem}
\begin{proof}
连接三边中点, 把 $T$ 分成四个相似三角形. 内边积分互相抵消, 四个边界积分的和等于原积分. 至少一个小三角形积分模不小于原来的四分之一. 反复选择, 得嵌套闭三角形 $T_n$, 直径 $d/2^n$, 周长 $L/2^n$, 且 $|I(T)|\le4^n|I(T_n)|$. 它们交于一点 $a$, 这是紧性和直径趋零的结果.

在 $a$ 复可微, 可写 $f(z)=f(a)+f'(a)(z-a)+\eta(z)(z-a)$, 其中 $\eta(z)\to0$. 常数和一次项有原函数, 它们绕三角形积分为零. 当 $n$ 大到 $T_n$ 上 $|\eta|<\eps$ 时, 积分估计给 $|I(T_n)|\le\eps dL/4^n$. 因而 $|I(T)|\le\eps dL$. 对任意 $\eps>0$ 成立, 所以积分为零.
\end{proof}
在圆盘内选择基点, 对每点沿直线积分定义 $F$. 三角形定理说明改变最后一段时路径积分不变, 因而前节的差商论证仍给 $F'=f$. 所以全纯函数在每个圆盘上都有原函数. 这一局部事实将替代原课程 Green 证明中暂时使用的连续偏导条件.

\originalsection{同伦与有孔区域}
两条同端点路径在 $\Omega$ 内同伦, 指存在连续 $H(s,t)$ 将一条变为另一条, 且端点固定. 闭路同伦可以允许共同起点随变形移动. 域称为单连通, 若每条闭路都能在域内连续收缩为一点.
\begin{theorem}\label{thm:homotopy}
全纯函数沿同伦的分段光滑路径的积分相同. 特别地, 单连通域上所有闭路积分为零, 每个全纯函数有原函数.
\end{theorem}
\begin{proof}
同伦像是域内的紧集, 可用有限个有原函数的圆盘覆盖. 该有限覆盖在紧同伦像上有正的 Lebesgue 数. 参数方形上 $H$ 一致连续; 存在足够细的矩形网格, 使每格的像包含在一个这样的圆盘内. 对网格边, 用连接端点的直线段替代其连续像, 必要时再细分, 使相关线段与相邻格顶点都在相应圆盘内. 一格四边的积分为零, 因其在有原函数的圆盘内. 求和后内部边抵消, 只剩外边.

外边对应的原路径各小段和直线替代段也在一个圆盘内, 原函数保证积分一致. 固定端点的两条端边为常路径, 贡献为零; 闭路同伦时两条端边相互反向, 贡献也抵消. 因而变形前后积分相同. 将闭路收缩为一点便得零积分, 再用前节等价性得到原函数.
\end{proof}
这个证明并不假定连续同伦的每条中间路径都可微; 网格与局部原函数只需要端点和最终两条可积路径. 因此不会把光滑性暗中施加给一般拓扑同伦.

若一个有界区域边界由外圈 $C_0$ 和互不相交的内圈 $C_1,\ldots,C_m$ 构成, 各圈为两两不相交的分段光滑简单闭曲线, 所有内圈位于外圈内部且各自内部不重叠, 各圈都暂按逆时针方向记录, 则正向区域边界为 $C_0-C_1-\cdots-C_m$. 当 $f$ 在区域闭包邻域全纯时,
\[
\int_{C_0}f\dd z=\sum_{j=1}^m\int_{C_j}f\dd z.
\]
对多边形边界, 剖分为有限三角形即得; 一般分段光滑边界用域内足够近的多边形替代, 局部原函数保证小段积分保持, 因而同样成立. 等价地可沿连接各内圈的割线打开区域, 两份割线反向抵消. 积分方向规则始终是沿边走时区域在左侧.

\begin{center}
\begin{tikzpicture}[>=Stealth,scale=.9]
\draw[thick] (0,0) ellipse (3 and 1.7);
\draw[thick] (-1,0) circle (.55); \draw[thick] (1,0) circle (.5);
\draw[->] (0,1.7)--(-.45,1.68); \node at (2.7,1.6) {$C_0$};
\draw[->] (-1,-.55)--(-1.3,-.47); \node at (-1,-.95) {$-C_1$};
\draw[->] (1,-.5)--(.72,-.42); \node at (1,-.92) {$-C_2$};
\node at (0,1) {$\Omega$};
\end{tikzpicture}
\end{center}
\begin{remark}
对于 $C^1$ 的 $u,v$, Green 定理也给 $\int f\dd z=\iint(-v_x-u_y)\dd x\dd y+\ii\iint(u_x-v_y)\dd x\dd y=0$. 这是原课程路线的直接验证. Goursat 论证补足了仅假定全纯时的正则性问题.
\end{remark}

\originalsection{习题与解答}
\begin{exercise}
计算从1到 $1+\ii$ 的直线段上 $\int\dd z/z$, 以及单位圆上 $\int\dd z/z$.
\end{exercise}
\begin{solution}
第一条路径在主对数域内, 故积分为 $\Log(1+\ii)-\Log1=\tfrac12\log2+\pi\ii/4$. 单位圆绕过主对数的割线, 不能用同一全局原函数, 直接参数化得 $2\pi\ii$.
\end{solution}
\begin{exercise}
证明连续 $f$ 若在一个圆盘内沿所有三角形边界积分为零, 则在圆盘内全纯.
\end{exercise}
\begin{solution}
以基点为首端, 定义沿直线段的积分 $F(z)$. 三角形条件允许将差值化为从 $z$ 到 $z+h$ 的直线段积分, 因而 $F'=f$. 第4章会证明全纯 $F$ 的导数仍全纯, 于是 $f$ 全纯. 这一命题称为 Morera 定理; 在使用它时须先建立 Cauchy 导数公式.
\end{solution}

\originalchapter{Cauchy 公式与模估计}
\sourceinfo{18.04 Topic 4 全部正文; 18.112 L10为主题背景. 交换极限与全域常数论证为编者补充}
\originalsection{积分公式}
\begin{theorem}[Cauchy 积分公式]\label{thm:cauchy}
设 $f$ 在闭圆盘 $\overline{D(a,R)}$ 的某开邻域全纯, $C$ 是其逆时针边界. 对任意 $w\in D(a,R)$,
\[
f(w)=\frac1{2\pi\ii}\int_C\frac{f(z)}{z-w}\dd z.
\]
公式对其他分段光滑简单闭曲线也成立, 前提是 $f$ 在曲线及整个内部的开邻域全纯, $w$ 位于内部.
\end{theorem}
\begin{proof}
在 $w$ 周围挖去半径 $\rho$ 的小圆盘. 核函数在剩余闭区域邻域全纯, 有孔边界定理将外圈积分化为小圆 $C_\rho$ 积分. 分解 $f(z)=f(w)+(f(z)-f(w))$. 第一项积分为 $2\pi\ii f(w)$. 第二项的模不超过 $2\pi\max_{|z-w|=\rho}|f(z)-f(w)|$, 由连续性趋零. 外圈积分不依赖 $\rho$, 所以等于第一项. 一般曲线同样挖小圆即可.
\end{proof}
要点是积分不只看边界上的全纯性; 内部若有奇点, 不可套用公式. 例如 $1/z$ 在单位圆附近全纯, 但单位圆内部有奇点, 因而其积分不为零.
\begin{theorem}[导数公式]\label{thm:derivcauchy}
同一条件下 $f$ 有各阶导数, 且
\[
f^{(n)}(w)=\frac{n!}{2\pi\ii}\int_C\frac{f(z)}{(z-w)^{n+1}}\dd z,\qquad n=0,1,2,\ldots.
\]
\end{theorem}
\begin{proof}
固定一个严格在圆内的紧集 $K$, 它到 $C$ 的距离为 $d>0$. 当 $|h|<d/2$ 时,
\[
\frac{1}{h}\left(\frac1{z-w-h}-\frac1{z-w}\right)=\frac1{(z-w-h)(z-w)}.
\]
右端在 $z\in C,w\in K$ 上一致趋于 $(z-w)^{-2}$, 差的模为 $O(|h|/d^3)$. 乘以有界 $f(z)$ 后积分差仍趋零, 故可把差商极限移入积分. 得到一阶公式. 同样的距离估计允许再次求导, 而核函数每次求导产生因子 $n+1$, 归纳即得各阶公式. 核导数连续且一致受界, 积分表示又表明各阶导数连续.
\end{proof}
复可微一次便自动无限次可微. 这也说明全纯函数的实、虚部有连续的任意阶偏导, 前章 Green 路线中的正则性最终获得保证, 但不作为该路线未经证明的起点.
\begin{example}
计算逆时针圆 $|z|=2$ 上 $\int\ee^z/(z-1)^3\dd z$.
\end{example}
\begin{solution}
分子整, 1在圆内. 用 $n=2$ 的公式, 积分为 $2\pi\ii f''(1)/2!=\pi\ii\ee$. 若圆改为 $|z|=1/2$, 被积函数在圆内全纯, 积分为0.
\end{solution}

\originalsection{Cauchy 估计与 Liouville 定理}
\begin{corollary}[Cauchy 估计]\label{cor:cauchyestimate}
若 $f$ 在 $\overline{D(a,R)}$ 邻域全纯, $M_R=\max_{|z-a|=R}|f(z)|$, 则 $|f^{(n)}(a)|\le n!M_R/R^n$.
\end{corollary}
\begin{proof}
导数公式中的分母模为 $R^{n+1}$, 圆长度为 $2\pi R$. 用积分估计即可.
\end{proof}
\begin{theorem}[Liouville]
有界整函数必为常数.
\end{theorem}
\begin{proof}
若 $|f|\le M$, 则以任意 $a$ 为心、任意 $R$ 为半径的圆都适用估计, 给 $|f'(a)|\le M/R$. 令 $R\to\infty$ 得 $f'(a)=0$. 前章的零导数结论给出常数性.
\end{proof}
\begin{corollary}[代数基本定理]
次数为 $n\ge1$ 的复多项式按重数恰有 $n$ 个复根.
\end{corollary}
\begin{proof}
若 $P$ 无根, 则 $1/P$ 整. 最高次项给 $|P(z)|\to\infty$ ($|z|\to\infty$), 所以 $1/P$ 在大圆外有界; 圆内由连续性和紧性也有界. Liouville 迫使 $1/P$ 恒定, 与 $P$ 次数正矛盾. 因而有一根 $a$, 多项式除法给 $P=(z-a)Q$. 反复对降次多项式应用存在根结论, 直到次数0; 这给出 $n$ 个一次因子, 同时也证明不可能有额外根.
\end{proof}
更一般地, 若 $m\in\mathbb Z_{\ge0}$, 整函数满足 $|f(z)|\le C(1+|z|^m)$, 取 $n>m$ 后 Cauchy 估计令 $R\to\infty$, 得所有这些阶的导数为零. 第5章的 Taylor 展开说明 $f$ 是次数至多 $m$ 的多项式.

\originalsection{平均值与最大模}
\begin{proposition}[平均值性质]
若 $f$ 在闭圆盘邻域全纯, 则 $f(a)=(2\pi)^{-1}\int_0^{2\pi}f(a+r\ee^{\ii t})\dd t$.
\end{proposition}
\begin{proof}
在 Cauchy 公式中代入 $z=a+r\ee^{\ii t}$, $\dd z=\ii r\ee^{\ii t}\dd t$, 核的分母恰好抵消.
\end{proof}
\begin{theorem}[最大模原理]\label{thm:maxmod}
域 $\Omega$ 上全纯函数的模若在域内某点有局部最大值, 则函数恒定. 若 $\Omega$ 有界且 $f$ 连续到闭包, 最大模在边界取得; 非常数时内部不会取得它.
\end{theorem}
\begin{proof}
设 $|f(a)|=M$ 为一个小圆盘内的最大模. $M=0$ 时此圆盘上 $f=0$. $M>0$ 时乘一个单位复数, 可设 $f(a)=M$ 为正实数. 平均值性质的实部给 $M$ 等于圆周上 $\Re f$ 的平均, 而 $\Re f\le |f|\le M$. 连续的非负差 $M-\Re f$ 平均为零, 故每点都为零; 再用 $|f|\le M$ 得 $\Im f=0$. 对所有足够小的圆成立, 所以邻域内 $f=M$.

局部常数如何传遍域将在第5章恒等定理给出完整证明. 也可现在作如下论证: 令 $S$ 为具有某常数值且在邻域恒定的点集. 它开; 若 $a$ 在域内为 $S$ 的极限点, 则 $f$ 的各阶导数在这些邻域为零, 连续性使它们在 $a$ 都为零. 第5章从 Cauchy 公式直接得到的局部 Taylor 展开使 $a$ 也属于 $S$. 所以 $S$ 又闭, 连通性给 $S=\Omega$. 这只依赖下一章已经由公式可证明的展开, 不反过来用最大模证明展开. 有界闭包紧, $|f|$ 取得最大值, 非常数时内部不能取得, 只剩边界.
\end{proof}
原理在无界域上的边界表述需要额外的无穷远控制. 例如 $\ee^{-\ii z}$ 在上半平面实轴上的模为1, 内部模为 $\ee^{\Im z}>1$, 整个域上并没有有限最大值. 最小模也不能直接照抄: $f(z)=z$ 在0取得模的最小值, 却非常数. 若 $f$ 无零点, 可对 $1/f$ 应用最大模原理.

\originalsection{Morera 与极限}
\begin{theorem}[Morera]\label{thm:morera}
开集上连续 $f$ 若对每个闭三角形及其内部包含于开集的三角形, 边界积分都为零, 则 $f$ 全纯.
\end{theorem}
\begin{proof}
在每个小圆盘内用三角形条件构造原函数 $F$, 第3章差商计算给 $F'=f$. 导数公式保证 $F'$ 全纯. 因而 $f$ 局部全纯, 也就全纯于原开集.
\end{proof}
\begin{theorem}\label{thm:weierstrasslimit}
开集上全纯函数列 $f_j$ 若在每个紧集上一致收敛于 $f$, 则 $f$ 全纯, 且对每个固定 $n$, $f_j^{(n)}$ 在紧集上一致收敛于 $f^{(n)}$.
\end{theorem}
\begin{proof}
局部一致极限连续, 且三角形积分可以取极限, 故 Morera 给 $f$ 全纯. 对一个圆盘内更小的紧集, 用同一外圆的导数公式. 差的积分模至多为 $n!\ell(C)\sup_C|f_j-f|/(2\pi d^{n+1})$, 其中 $d$ 是紧集到圆的距离. 故导数一致收敛. 任意紧集用有限个这样的小圆盘覆盖便得结论.
\end{proof}
\originalsection{习题与解答}
\begin{exercise}
整函数满足 $\Re f\le A$. 证明 $f$ 恒定.
\end{exercise}
\begin{solution}
$\ee^f$ 整且模不超过 $\ee^A$, 所以恒定. 其导数 $\ee^f f'=0$, 指数无零点, 故 $f'=0$.
\end{solution}
\begin{exercise}
若 $f$ 在闭单位圆盘邻域全纯且 $|f|\le1$ 于单位圆周, 证明 $|f'(0)|\le1$.
\end{exercise}
\begin{solution}
Cauchy 估计直接给 $|f'(0)|\le1$. 条件若只是在圆内某条线段成立, 则不能替代圆周条件.
\end{solution}
\begin{exercise}
证明整函数若 $f(z)\to0$ ($|z|\to\infty$), 则 $f\equiv0$.
\end{exercise}
\begin{solution}
大圆外有界, 圆内紧性给有界, 所以整函数恒定. 其无穷远极限使常数为0. 也可固定 $a$, 用大圆 Cauchy 公式估计或最大模原理令边界最大模趋零.
\end{solution}

\originalchapter{幂级数与恒等定理}
\sourceinfo{18.04 Topic 7, pp.1--10及未完成的收敛附录; 18.112 L4. 一致收敛估计为编者补充}
\originalsection{收敛半径}
\begin{definition}
级数 $\sum c_n$ 收敛, 指部分和趋于有限复数. 绝对收敛指 $\sum|c_n|<\infty$. 函数级数在集合 $K$ 上一致收敛, 指尾和可由同一个与点无关的界控制. 在每个紧集上一致收敛称为局部一致收敛.
\end{definition}
绝对收敛保证 Cauchy 条件: 任何足够后的尾和模不超过绝对值尾和. 复数完备性来自两个实坐标的完备性. 若 $|f_n(z)|\le M_n$ 于 $K$ 且 $\sum M_n<\infty$, 则尾和一致趋零; 这是 Weierstrass 判别法. 一致收敛使连续性和有限长度曲线上的积分可以交换极限, 因为积分误差至多为长度乘上一致误差.

比值判别中若 $|c_{n+1}/c_n|\to q<1$, 则后续项由几何级数控制, 绝对收敛; 若 $q>1$, 项不趋零, 级数发散. 根值判别使用 $\limsup |c_n|^{1/n}$, 无需极限存在. 值等于1时两者都不作结论.
\begin{theorem}[收敛半径]\label{thm:radius}
设 $\sum_{n\ge0}a_n(z-a)^n$ 为幂级数, 令 $R^{-1}=\limsup_{n\to\infty}|a_n|^{1/n}$, 按 $1/0=\infty$, $1/\infty=0$ 解释. 级数在 $|z-a|<R$ 绝对且局部一致收敛, 在 $|z-a|>R$ 发散. $R>0$ 时其和全纯, 可逐项求导、积分, 导数级数仍有半径 $R$.
\end{theorem}
\begin{proof}
对 $r<R$, 选 $r<\rho<R$. 由上极限, 对充分大 $n$ 有 $|a_n|\le \rho^{-n}$. 在 $|z-a|\le r$ 上项由 $(r/\rho)^n$ 控制, 所以一致绝对收敛. 若 $|z-a|>R$, 选择 $R<\rho<|z-a|$, 则有无穷多个 $n$ 满足 $|a_n|>\rho^{-n}$, 对应项模不趋零. 端点 $R=0,\infty$ 按同一论证处理.

部分和都是全纯多项式, 局部一致极限定理给和全纯、导数收敛. 导数系数 $na_n$ 的 $n$ 次根多一个趋于1的因子, 移动一个指标不改变半径; 也可用 $n(r/\rho)^n$ 的可求和性直接得到导数一致收敛. 对任意固定分段光滑路径, 其像是收敛圆盘内的紧集, 积分误差估计允许逐项积分.
\end{proof}
边界需要另查. $\sum z^n$ 在单位圆周处处发散; $\sum z^n/n^2$ 在闭单位圆盘上一致绝对收敛. 两者半径均为1. 不可把圆内局部一致收敛误说成在整个开圆盘上一致收敛.

\originalsection{Taylor 展开}
\begin{theorem}\label{thm:taylor}
若 $f$ 在 $D(a,R)$ 全纯, 则
\[
f(z)=\sum_{n=0}^\infty\frac{f^{(n)}(a)}{n!}(z-a)^n,\qquad |z-a|<R.
\]
展开在每个较小闭圆盘上一致绝对收敛, 且系数唯一.
\end{theorem}
\begin{proof}
取 $r<\rho<R$, 对 $|z-a|\le r$, $|\zeta-a|=\rho$, 几何级数给
\[
\frac1{\zeta-z}=\sum_{n=0}^N\frac{(z-a)^n}{(\zeta-a)^{n+1}}
 +\frac{(z-a)^{N+1}}{(\zeta-a)^{N+1}(\zeta-z)}.
\]
代入 Cauchy 公式, 系数积分由导数公式等于 $f^{(n)}(a)/n!$. 余项模不超过 $M_\rho(r/\rho)^{N+1}\rho/(\rho-r)$, 一致趋零. 绝对收敛由 Cauchy 系数估计得到. 任意幂级数表达可以逐项求导并在 $a$ 取值, 必须具有上述系数, 所以唯一.
\end{proof}
该证明只依赖 Cauchy 公式与几何级数, 并未使用最大模原理. 它补足第4章从局部常数到整个域恒定的最后环节, 因而不存在循环依赖. 全纯与解析的等价也至此建立: Taylor 定理给一个方向, 幂级数和的全纯性给反方向.

几个常用展开为
\begin{align*}
\ee^z&=\sum_{n\ge0}\frac{z^n}{n!},&
\sin z&=\sum_{n\ge0}\frac{(-1)^nz^{2n+1}}{(2n+1)!},\\
\cos z&=\sum_{n\ge0}\frac{(-1)^nz^{2n}}{(2n)!},&
\Log(1+z)&=\sum_{n\ge1}\frac{(-1)^{n-1}z^n}{n},\quad |z|<1.
\end{align*}
最后一式可从 $1/(1+z)$ 的几何级数逐项积分, 用 $\Log1=0$ 确定常数. 所使用的分支在整个圆盘内全纯. 展开中心改变时必须重新写成 $z-a$ 的幂; 如 $1/(1-z)$ 在 $a=5$ 的展开为 $-\tfrac14\sum_{n\ge0}(-(z-5)/4)^n$, 半径4.

\begin{example}
求 $\ee^z/(1-z)$ 在0的 Taylor 系数和半径.
\end{example}
\begin{solution}
在 $|z|<1$ 内两个级数绝对收敛, 其乘积可按总次数收集. 系数为 $a_n=\sum_{k=0}^n1/k!$. 因此展开是 $\sum a_nz^n$. $a_n\to\ee\ne0$, 根值判别给半径1. 从函数看, 点1是不能消去的极点; 若半径大于1, 幂级数和会在1附近全纯且与原函数在交域相同, 矛盾.
\end{solution}

\originalsection{零点与恒等}
\begin{definition}
若 $f(a)=0$ 且存在最小正整数 $m$ 使 $f^{(m)}(a)\ne0$, 则零点阶数为 $m$. 此时 $f(z)=(z-a)^mg(z)$, $g$ 在 $a$ 附近全纯, $g(a)\ne0$.
\end{definition}
\begin{theorem}[恒等定理]\label{thm:identity}
域上的全纯函数若零点在域内有聚点, 则恒为零. 因而非常零全纯函数的零点都孤立. 两个全纯函数若在域内有聚点的集合上相等, 则在整个域相等.
\end{theorem}
\begin{proof}
在零点聚点 $a$ 的 Taylor 展开若有首个非零系数, 则可分解为 $(z-a)^mg(z)$, 且小邻域内 $g\ne0$, 所以 $a$ 是孤立零点, 矛盾. 因而全部系数为零, $f$ 在 $a$ 邻域恒为零.

令 $S$ 为 $f$ 在邻域恒为零的点集. $S$ 开且非空. 若 $b$ 是 $S$ 在域内的极限点, 所有阶导数在 $S$ 为零, 由连续性在 $b$ 也为零, Taylor 展开使 $b\in S$. 所以 $S$ 在域内又闭, 连通性迫使它是全域. 最后应用于两函数之差.
\end{proof}
聚点必须位于域内. $\sin(\pi/z)$ 在穿孔平面全纯, 零点 $1/n$ 聚到缺失的0, 并不恒为零. 连通性也不能删去: 在上下半平面分别定义不同常数, 可使两个全纯函数只在一个分支相等.

若 $f$ 在0的全部导数为零, 则全纯性迫使它局部恒为零. 这与实光滑函数不同: 实函数 $\ee^{-1/x^2}$ 在 $x\ne0$ 定义、在0补零后无限次可微, 零点所有导数为零却不恒为零. 决定复函数的刚性的是复可微和 Cauchy 理论, 不是仅仅导数的数量.

\originalsection{幂级数解微分方程}
解析系数微分方程可借未知 Taylor 系数的递推求解, 但形式递推后还须验证收敛. 以 $f''+f=0$, $f(0)=A,f'(0)=B$ 为例, 写 $f=\sum a_nz^n$, 逐次系数给 $(n+2)(n+1)a_{n+2}+a_n=0$. 初值使 $a_0=A,a_1=B$, 因而
\[
a_{2n}=\frac{(-1)^nA}{(2n)!},\qquad
 a_{2n+1}=\frac{(-1)^nB}{(2n+1)!}.
\]
所得两列半径无穷, 逐项微分合法, 所以 $f=A\cos z+B\sin z$ 确实是整解. 任意全纯解都有 Taylor 展开并满足同一递推, 所以在0附近唯一, 恒等定理使其在连通共同域唯一. 若方程在某点的最高导数系数为零, 该点递推可能失效, 不能把这个正常点例子的结论无条件推广到奇点.

例如 $f'=(1-z)^{-1}f$, $f(0)=1$, 在 $|z|<1$ 有系数关系 $(n+1)a_{n+1}=\sum_{k=0}^na_k$. 由归纳 $a_n=1$, 和为 $1/(1-z)$, 在0的收敛半径1恰受方程系数的奇点限制. 形式系数并不保证整平面收敛, 需像这里一样另作半径验证.

\originalsection{习题与解答}
\begin{exercise}
求 $\sum_{n\ge0}n!z^n$ 和 $\sum_{n\ge0}z^n/(n!)^2$ 的收敛半径.
\end{exercise}
\begin{solution}
第一列相邻项模比为 $(n+1)|z|$, 任意 $z\ne0$ 时最终大于1, 所以半径0. 第二列比值为 $|z|/(n+1)^2\to0$, 任意 $z$ 都绝对收敛, 半径无穷.
\end{solution}
\begin{exercise}
证明整函数若在所有实数上等于 $\cos x$, 则等于 $\cos z$ 于全平面; 若仅在所有整数上等于 $\cos n$, 能否得到同一结论?
\end{exercise}
\begin{solution}
实轴具有复平面内的聚点, 恒等定理适用. 整数集合没有有限聚点, $\cos z+\sin(\pi z)$ 是反例.
\end{solution}
\begin{exercise}
求 $(\sin z-z)/z^3$ 在0的连续补值与前两项 Taylor 展开.
\end{exercise}
\begin{solution}
正弦级数给函数为 $-1/6+z^2/120+O(z^4)$. 补值为 $-1/6$, 补后在整个平面全纯; 除法造成的零分母已经被三阶零点消去.
\end{solution}

\originalchapter{Laurent 级数与留数}
\sourceinfo{18.04 Topic 7, pp.10--18及 Topic 8 全部正文; 18.112 L11、L14. 奇点分类判据证明为编者补充}
\originalsection{环域展开}
\begin{theorem}[Laurent]\label{thm:laurent}
若 $f$ 在环域 $r<|z-a|<R$ 全纯, 则有唯一表达 $f(z)=\sum_{n\in\Z}c_n(z-a)^n$, 正负两部分分别局部一致绝对收敛. 系数为
\[
c_n=\frac1{2\pi\ii}\int_{|\zeta-a|=\rho}\frac{f(\zeta)}{(\zeta-a)^{n+1}}\dd\zeta,\qquad r<\rho<R.
\]
圆逆时针, 系数与 $\rho$ 无关.
\end{theorem}
\begin{proof}
选 $r<\rho_1<|z-a|<\rho_2<R$. 对两圈间区域在 $z$ 挖小孔, Cauchy 公式给 $f(z)$ 等于外圈核积分减去内圈核积分. 外圈上按 $(z-a)/(\zeta-a)$ 展开几何级数, 得非负幂. 内圈上用
\[
\frac1{\zeta-z}=-\sum_{k\ge0}\frac{(\zeta-a)^k}{(z-a)^{k+1}},
\]
得到负幂; 负号与内圈前的负号相抵消. 两种展开在严格处于两圈间的紧子环域上一致绝对收敛, 所以可逐项积分. 系数公式的圈半径可变, 因被积函数在两圈间全纯且边界积分相同. 任意紧子环域均可这样选圈, 所以展开覆盖整个环域.

若有另一展开, 在环域内任一圈上乘 $(z-a)^{-m-1}$ 再积分, 一致收敛允许逐项积分. 整数幂圆积分只有指数 $-1$ 不为零, 所以必得到 $2\pi\ii c_m$, 证明唯一性.
\end{proof}
负幂部分叫主部, 非负幂部分叫正则部. Laurent 展开不仅依赖中心, 也依赖所用环域. 例如
\[
\frac1{z(z-1)}=-\frac1z\frac1{1-z}=-\sum_{n\ge0}z^{n-1},\quad 0<|z|<1,
\]
而在 $|z|>1$ 它等于 $\sum_{n\ge0}z^{-n-2}$. 外环的 $z^{-1}$ 系数为0, 不能当作0处孤立奇点的留数: 外圈已经同时绕住1处奇点.

\originalsection{孤立奇点}
\begin{definition}
若 $f$ 在 $0<|z-a|<R$ 全纯, 则 $a$ 是孤立奇点. 主部全为零称为可去奇点; 主部有限且最高负幂为 $(z-a)^{-m}$, 系数非零, 称为 $m$ 阶极点; 主部有无穷多非零项称为本性奇点.
\end{definition}
\begin{theorem}\label{thm:removable}
孤立奇点可去当且仅当 $f$ 在某穿孔邻域有界. $a$ 为 $m$ 阶极点当且仅当 $f(z)=(z-a)^{-m}g(z)$, $g$ 在 $a$ 附近全纯且 $g(a)\ne0$. 孤立奇点是极点当且仅当 $|f(z)|\to\infty$ ($z\to a$).
\end{theorem}
\begin{proof}
可去时补后的全纯函数连续, 故局部有界. 反向设 $|f|\le M$. 对 $k\ge1$, 负系数公式给 $|c_{-k}|\le M\rho^k$. 令 $\rho\downarrow0$ 得所有负系数为零, 正则部提供全纯补值. 极点分解直接由有限主部得到; 逆向把 $g$ 作 Taylor 展开即可, 它的常数项非零保证阶数恰为 $m$.

极点分解给 $|f|\to\infty$. 反向若模趋无穷, 则足够近处 $f\ne0$, $1/f$ 全纯且趋零, 可补为全纯零点. 它不能局部恒为零, 所以有有限正阶 $m$, 分解 $1/f=(z-a)^mh$ 且 $h(a)\ne0$, 得极点.
\end{proof}
函数 $\sin z/z$ 在0可去, $1/z^m$ 是 $m$ 阶极点, $\ee^{1/z}$ 是本性奇点. 本性并不意味着 ``模总趋向无穷'': 沿正实方向 $\ee^{1/z}\to\infty$, 沿负实方向趋零. $\Log z$ 的0点则不是孤立奇点, 因为没有覆盖整个穿孔圆盘的单值全纯分支.
\begin{theorem}[Casorati--Weierstrass]
若 $a$ 为本性奇点, 则任意穿孔邻域内 $f$ 的值在 $\C$ 中稠密.
\end{theorem}
\begin{proof}
若值域不稠密, 存在 $b$ 和 $\eps>0$, 使整个小穿孔邻域上 $|f-b|\ge\eps$. 则 $g=1/(f-b)$ 有界全纯, 可补过 $a$. 若 $g(a)\ne0$, 则 $f=b+1/g$ 可去; 若 $g(a)=0$, $g$ 有有限阶零点, 则 $f$ 为极点. 两者均与本性矛盾.
\end{proof}
无穷远奇点通过 $f(1/\zeta)$ 分类. 非常数多项式在无穷远为极点, 有界于无穷远邻域的整函数为常数, 非多项式整函数在无穷远为本性奇点.

\originalsection{留数与计算}
\begin{definition}
孤立奇点 $a$ 的留数是穿孔圆盘 Laurent 展开的 $c_{-1}$, 记 $\Res(f,a)$. 等价地, 在逆时针小圆上有
\[
\Res(f,a)=\frac1{2\pi\ii}\int_{|z-a|=\rho}f(z)\dd z.
\]
\end{definition}
在 $m$ 阶极点, 写 $g(z)=(z-a)^mf(z)$, 则 $\Res(f,a)=g^{(m-1)}(a)/(m-1)!$. 因为所需系数恰是 $g$ 的 $m-1$ 次项. 若 $f=p/q$, $q(a)=0$, $q'(a)\ne0$, 则 $\Res(f,a)=p(a)/q'(a)$; 即使 $p(a)=0$ 导致可去, 该式仍给零留数.
\begin{example}
求 $\ee^z/z^3$, $\cot z$, $\ee^{1/z}$ 在0的留数.
\end{example}
\begin{solution}
第一式的留数为指数级数的 $z^2$ 系数 $1/2$. $\cot z=\cos z/\sin z$, 分母在0有简单零点, 所以留数为1. 本性奇点的留数仍有意义: $\ee^{1/z}=1+z^{-1}+z^{-2}/2!+\cdots$, 留数也为1. 极点阶数与留数是不同信息, 留数为零也不代表可去, 如 $1/z^2$.
\end{solution}
长除法还给 $\cot z=z^{-1}-z/3-z^3/45+O(z^5)$. 可把未知 Laurent 系数乘回 $\sin z$ 并与 $\cos z$ 的系数比较, 不需反复求高阶导数.

\originalsection{留数定理}
\begin{theorem}\label{thm:residue}
设 $C$ 为分段光滑简单闭曲线, 正向绕行. 若 $f$ 在 $C$ 及其内部的开邻域除有限个孤立奇点 $a_1,\ldots,a_m$ 外全纯, 且奇点均不在 $C$ 上, 则 $\int_Cf\dd z=2\pi\ii\sum_j\Res(f,a_j)$.
\end{theorem}
\begin{proof}
围每个奇点选互不相交的小圆, 完全在曲线内, 且各自穿孔圆盘内无其他奇点. 在挖孔后的区域上应用有孔边界定理, 外圈积分等于各小圈的逆时针积分之和. 每个小圈积分由留数定义等于 $2\pi\ii\Res(f,a_j)$.
\end{proof}
奇点允许是本性奇点. 若曲线自交或多次绕行, 应按绕数加权, 不能只说 ``曲线里面''. 第8章给出明确的绕数定义与公式.

\begin{example}
求 $\int_{|z|=2}(5z-2)/(z(z-1))\dd z$ 和 $\int_{|z|=1}z^2\sin(1/z)\dd z$, 均逆时针.
\end{example}
\begin{solution}
第一式两极点0、1的留数分别为2、3, 所以积分为 $10\pi\ii$. 第二式在0的展开为 $z-1/(6z)+1/(120z^3)-\cdots$, 留数为 $-1/6$, 积分为 $-\pi\ii/3$.
\end{solution}
无穷远留数处理的是微分 $f(z)\dd z$, 不只是函数. 在坐标 $\zeta=1/z$ 下 $\dd z=-\zeta^{-2}\dd\zeta$, 定义 $\Res(f,\infty)=-\Res(\zeta^{-2}f(1/\zeta),0)$. 若平面只有有限个孤立奇点, 大圆的 Laurent 展开给 $\Res(f,\infty)$ 为 $z^{-1}$ 系数的相反数. 因而 $\sum_{a\in\C}\Res(f,a)+\Res(f,\infty)=0$. 这里负号同时来自坐标微分和内外区域的方向变换.

\originalsection{零极阶数与本性值域}
若全纯 $p,q$ 在 $a$ 的零点阶数分别为 $m,n$, 则写 $p=(z-a)^mu$, $q=(z-a)^nv$, $u(a)v(a)\ne0$, 得 $p/q=(z-a)^{m-n}u/v$. 因而 $m>n$ 时有 $m-n$ 阶零点, $m=n$ 时可去且补值非零, $m<n$ 时有 $n-m$ 阶极点. 这条一般规则统一了有理函数中约因子与零分母的关系.

原课还陈述大 Picard 定理: 在本性奇点的任意穿孔邻域内, 函数取每个复数值无穷多次, 至多允许一个例外值. 它比 Casorati 的稠密性强得多, 例如 $\ee^{1/z}$ 恰遗漏0. 这里保留原课的引用版陈述; 其证明需要超出本主课的值分布工具, 不作为本书已证明的核心定理, 后文也不依赖它. 本书已完整证明并实际使用的是 Casorati 定理, 不把稠密性冒称为 Picard.

\originalsection{习题与解答}
\begin{exercise}
分类 $(\ee^z-1)/z^2$, $\sin(1/z)$ 和 $z\sin(1/z)$ 在0的奇点并计算留数.
\end{exercise}
\begin{solution}
第一式展开为 $z^{-1}+1/2+\cdots$, 简单极点, 留数1. 第二式主部无穷, 本性, 留数1. 第三式为 $1-1/(6z^2)+\cdots$, 仍本性, 留数0.
\end{solution}
\begin{exercise}
求 $1/[z(z^2+1)(z-2)^2]$ 在2的留数.
\end{exercise}
\begin{solution}
二阶极点处 $g=1/[z(z^2+1)]$. $g'=- (3z^2+1)/[z^2(z^2+1)^2]$, 所以留数 $g'(2)=-13/100$.
\end{solution}
\begin{exercise}
证明整函数在无穷远为可去奇点则为常数, 为 $m$ 阶极点则为 $m$ 次多项式.
\end{exercise}
\begin{solution}
可去意味着大圆外有界, 全平面有界后用 Liouville. $m$ 阶极点给 $|f(z)|\le C|z|^m$ 于大圆外, 由 Cauchy 估计消去所有高于 $m$ 的 Taylor 系数. 主部最高阶非零又使次数恰为 $m$.
\end{solution}

\originalchapter{留数的积分应用}
\sourceinfo{18.04 Topic 9 全部正文; 18.112 L15. 各类大、小圆弧估计与 Fourier 反演条件为编者补充}
\originalsection{实轴积分}
把实积分嵌入闭曲线, 留数定理负责计算闭路积分, 另一个同样必要的工作是证明增加的边界积分趋零. 如果只写 ``闭合曲线后取留数'', 常会漏掉收敛条件、方向或实轴上的奇点.
\begin{proposition}
有理函数 $R=P/Q$ 无实极点, 且 $\deg Q\ge\deg P+2$. 则
\[
\int_{-\infty}^\infty R(x)\dd x=2\pi\ii\sum_{\Im a>0}\Res(R,a).
\]
\end{proposition}
\begin{proof}
实轴上 $R=O(x^{-2})$, 故绝对可积. 在上半圆闭路上应用留数定理. 对足够大的半径, 圆弧上 $|R|\le C/R^2$, 弧长 $\pi R$, 积分模不超过 $C\pi/R\to0$. 选半径超过所有极点后, 留数和固定, 实线段积分趋于整条实轴积分.
\end{proof}
\begin{example}
计算 $\int_{-\infty}^{\infty}(x^2+1)^{-2}\dd x$ 和 $\int_0^\infty (x^4+1)^{-1}\dd x$.
\end{example}
\begin{solution}
第一式上半平面只有二阶极点 $\ii$. $g(z)=(z+\ii)^{-2}$, $g'(\ii)=-2/(2\ii)^3=1/(4\ii)$, 积分为 $2\pi\ii/(4\ii)=\pi/2$.

第二式的上半平面极点为 $a=\ee^{\ii\pi/4}$ 和 $b=\ee^{3\ii\pi/4}$. 留数为 $1/(4a^3)$、$1/(4b^3)$, 和为 $-\ii/(2\sqrt2)$. 整条实轴积分为 $\pi/\sqrt2$, 被积函数偶, 所以半轴积分为 $\pi/(2\sqrt2)$.
\end{solution}

\originalsection{振荡因子与 Jordan 估计}
\begin{lemma}[Jordan 估计]\label{lem:jordan}
固定 $t>0$. 若上半圆 $C_R$ 上 $|g(z)|\le M_R$ 且 $M_R\to0$, 则 $\int_{C_R}\ee^{\ii tz}g(z)\dd z\to0$.
\end{lemma}
\begin{proof}
设 $z=R\ee^{\ii\theta}$, 得积分模不超过 $M_RR\int_0^\pi\ee^{-tR\sin\theta}\dd\theta$. 在 $0\le\theta\le\pi/2$ 有 $\sin\theta\ge2\theta/\pi$, 利用对称性, 上界为
\[
2M_RR\int_0^{\pi/2}\ee^{-2tR\theta/\pi}\dd\theta
\le\frac{\pi M_R}{t}\longrightarrow0.
\]
\end{proof}
$t<0$ 要在下半平面闭合, 因为那里的指数模衰减; 闭路为顺时针, 需要负号. 此估计允许 $g=O(1/z)$, 比纯粹的长度乘最大值估计更强.
\begin{example}
对 $a>0$, 计算 $\int_{-\infty}^\infty\ee^{\ii tx}/(x^2+a^2)\dd x$.
\end{example}
\begin{solution}
$t>0$ 时上半平面留数为 $\ee^{-at}/(2\ii a)$, 结果为 $\pi\ee^{-at}/a$. $t<0$ 时下半平面留数为 $\ee^{at}/(-2\ii a)$, 再乘顺时针的 $-2\pi\ii$, 得 $\pi\ee^{at}/a$. $t=0$ 直接用有理积分. 故统一结果为 $\pi\ee^{-a|t|}/a$, 实部给余弦积分, 虚部因奇偶性为0.
\end{solution}

\originalsection{圆周换元与割线积分}
对于三角有理积分, 令 $z=\ee^{\ii\theta}$, 则 $\dd\theta=\dd z/(\ii z)$, $\cos\theta=(z+z^{-1})/2$, $\sin\theta=(z-z^{-1})/(2\ii)$. 原积分化为单位圆积分, 但要连同换元产生的 $z$ 因子一起约分再找奇点.
\begin{example}
设 $a,b\in\R$, $a>|b|$, 计算 $I=\int_0^{2\pi}\dd\theta/(a+b\cos\theta)$.
\end{example}
\begin{solution}
$b=0$ 时结果为 $2\pi/a$. $b\ne0$ 时换元得 $I=(2/\ii)\int_{|z|=1}\dd z/(bz^2+2az+b)$. 两根为 $(-a\pm\sqrt{a^2-b^2})/b$, 乘积1, 只有加号根在圆内. 在该根分母导数为 $2\sqrt{a^2-b^2}$, 故 $I=2\pi/\sqrt{a^2-b^2}$.
\end{solution}

\begin{example}\label{ex:keyhole}
对 $0<\alpha<1$, 计算 $I_\alpha=\int_0^\infty x^{\alpha-1}/(1+x)\dd x$.
\end{example}
\begin{solution}
取辐角 $0<\arg z<2\pi$, 使割线位于正实轴. 在割平面上积分 $z^{\alpha-1}/(1+z)$, 沿钥匙孔曲线: 上岸从 $\eps$ 到 $R$, 大圆逆时针, 下岸从 $R$ 返回 $\eps$, 小圆顺时针. 上岸值趋 $x^{\alpha-1}/(1+x)$, 下岸多因子 $\ee^{2\pi\ii(\alpha-1)}=\ee^{2\pi\ii\alpha}$. 两岸贡献之和为 $(1-\ee^{2\pi\ii\alpha})I_\alpha$.

大圆积分为 $O(R^{\alpha-1})\to0$, 小圆为 $O(\eps^\alpha)\to0$, 这恰好用到 $0<\alpha<1$. 唯一极点 $-1$ 的留数为 $\ee^{\pi\ii(\alpha-1)}=-\ee^{\pi\ii\alpha}$. 因而
\[
(1-\ee^{2\pi\ii\alpha})I_\alpha=-2\pi\ii\ee^{\pi\ii\alpha},\qquad
I_\alpha=\frac\pi{\sin\pi\alpha}.
\]
严格地说两岸可先取夹角 $\delta$ 的射线, 在固定 $\eps,R$ 下让 $\delta\to0$, 再令 $R\to\infty,\eps\to0$, 由绝对积分估计控制这些极限.
\end{solution}
\begin{center}
\begin{tikzpicture}[>=Stealth,scale=.75]
\draw[->] (-3.8,0)--(4.5,0) node[right]{$\Re z$};
\draw[->] (0,-3.5)--(0,3.7) node[above]{$\Im z$};
\draw[thick] (.55,.13)--(3.5,.13) arc (2.13:357.87:3.5024)--(.55,-.13);
\draw[thick] (.55,-.13) arc (-13.3:-346.7:.565);
\draw[->,thick] (1.3,.13)--(2,.13); \draw[->,thick] (2,-.13)--(1.3,-.13);
\fill (-1,0) circle (2pt) node[above left]{$-1$};
\node at (2.15,.52){$\arg z\to0$}; \node at (2.2,-.57){$\arg z\to2\pi$};
\node at (-2.5,2.6){$|z|=R$}; \node at (.6,.8){$\eps$};
\end{tikzpicture}
\end{center}

\originalsection{实轴极点与主值}
定义 Cauchy 主值 $\operatorname{PV}\int_{-\infty}^\infty f(x)\dd x$ 为在实轴奇点两侧对称删去小区间并在无穷远对称截断后的极限. 它与通常的广义积分不同. 若 $f(z)=c/(z-a)+h(z)$, 在实轴点 $a$ 上方用顺时针小半圆从 $a-\eps$ 到 $a+\eps$ 避开极点, 则奇异部分积分为 $-\pi\ii c$, 正则部分积分趋零. 因而这一半留数贡献不能遗漏.
\begin{example}
证明 $\int_0^\infty\sin x/x\dd x=\pi/2$, 广义积分按通常单侧极限理解.
\end{example}
\begin{solution}
对 $\ee^{\ii z}/z$ 取上半圆和避开0的上方小半圆, 闭路内无奇点. 大圆贡献由 Jordan 估计趋零, 小半圆贡献趋 $-\pi\ii$, 故实轴主值为 $\pi\ii$. 虚部 $\sin x/x$ 在0可连续补值, 是偶函数. 其无穷远单侧积分收敛: 分部积分给 $\int_A^B\sin x/x\dd x=[-\cos x/x]_A^B-\int_A^B\cos x/x^2\dd x$, 尾量为 $O(1/A)$. 因此实轴虚部积分等于两倍半轴积分, 得所求值.
\end{solution}

\originalsection{级数与 Fourier 变换}
\begin{example}
证明对 $a>0$, $\sum_{n\in\Z}(n^2+a^2)^{-1}=\pi\coth(\pi a)/a$.
\end{example}
\begin{solution}
积分 $F(z)=\pi\cot(\pi z)/(z^2+a^2)$ 于顶点 $\pm(N+1/2)\pm\ii(N+1/2)$ 的方形, $N>a$. 整数极点处留数为 $1/(n^2+a^2)$. 另外两点 $\pm\ii a$ 的留数之和为 $-\pi\coth(\pi a)/a$, 由 $\cot(\ii t)=-\ii\coth t$ 得到.

在竖边 $\Re z=N+1/2$ 上, $|\cot\pi z|=|\tanh(\pi\Im z)|\le1$; 横边上用指数表达可得统一有界, 例如当 $|\Im z|\ge1/2$ 时不超过 $\coth(\pi/2)$. 分母在四边模至少为 $|z|^2-a^2\ge(N+1/2)^2-a^2$, 总边长为 $8(N+1/2)$. 积分趋零. 留数定理后取 $N\to\infty$, 整数项绝对收敛, 即得结果.
\end{solution}

取 Fourier 约定 $\widehat g(\omega)=\int_\R g(x)\ee^{-\ii\omega x}\dd x$. 前面振荡积分给 $g(x)=1/(x^2+a^2)$ 的变换 $\widehat g(\omega)=\pi\ee^{-a|\omega|}/a$. 反演中的常数取决于这一约定, 不能省略 $1/(2\pi)$.
\begin{theorem}[Gaussian 正则化反演]\label{thm:fourier}
若 $g\in L^1(\R)$, 在 $x$ 连续, 则
\[
g(x)=\lim_{\eta\downarrow0}\frac1{2\pi}\int_\R\widehat g(\omega)\ee^{\ii\omega x-\eta\omega^2}\dd\omega.
\]
若另有 $\widehat g\in L^1$, 可去掉正则化并直接积分. 分段光滑跳点的正则化值是左右极限平均.
\end{theorem}
\begin{proof}
实 Gaussian 积分及配方的 Fourier 计算给
\[
\frac1{2\pi}\int_\R\ee^{\ii\omega s-\eta\omega^2}\dd\omega
 =\frac1{\sqrt{4\pi\eta}}\ee^{-s^2/(4\eta)}=:K_\eta(s).
\]
该计算也可先对 $s$ 求导并分部积分得到微分方程 $I'(s)=-sI(s)/(2\eta)$, 用 $I(0)=\sqrt{\pi/\eta}$ 定常数. 因 $g$ 可积且 Gaussian 频率因子可积, 绝对积分允许换序, 正则化积分等于 $\int g(t)K_\eta(x-t)\dd t$.

$K_\eta\ge0$, 总积分1, 且任意 $\delta>0$ 外的质量趋零. 在 $|t-x|<\delta$ 用连续性令 $|g(t)-g(x)|$ 小; 外部积分中的 $g$ 项由 $\|g\|_1\sup_{|s|\ge\delta}K_\eta(s)\to0$ 控制, 常数项由尾质量控制. 故极限为 $g(x)$. 跳点时用核的偶性分为两半, 每半质量为 $1/2$, 同理得左右平均. 若 $\widehat g$ 可积, 被积函数由 $|\widehat g|$ 控制, 可取极限去正则化.
\end{proof}
这里 $L^1$ 表示绝对可积函数, 可先以本书例子的分段连续广义积分理解. 本证明调用绝对可积时的 Fubini 换序与控制收敛: 有可积绝对值上界时允许换序或取极限. 对这些具体核, 相同结论也可先在有界矩形上交换, 再以给出的尾估计取极限. 原课程对一般反演略去的收敛细节由这个版本补足; 不把只连续且指数增长的函数未经条件检查便代入一个普通绝对积分.

\originalsection{两个分支点的割线}
\begin{example}
计算 $\int_1^\infty\dd x/[x\sqrt{x^2-1}]$, 其中根为正实根.
\end{example}
\begin{solution}
在 $\C\setminus[-1,1]$ 选择 $S(z)=\sqrt{z^2-1}$, 满足 $S(z)\sim z$ 于无穷远. 可用 $S(z)=(z+1)\sqrt{(z-1)/(z+1)}$ 定义: 分式在割线外避开负实轴, 外根取主分支. $z>1$ 时 $S$ 为正, $z<-1$ 时为负. 考虑 $f=1/(zS(z))$, 它在割线外全纯且为 $O(z^{-2})$.

在割线端点 $\pm1$ 附近, $f$ 有平方根型增长; 0也位于割线上, 绕割线做闭路时还要处理该点. 采用该分支导出的原函数可以直接完成计算:
\[
\frac{\dd}{\dd z}\arccos(1/z)=\frac1{zS(z)},\qquad z>1,
\]
其中实区间上逆余弦从0连续增加到 $\pi/2$. 该导数由局部反三角函数的隐式微分得到, 沿实区间所选 $S>0$, 所以 $\int_{1+\eps}^R f(x)\dd x=\arccos(1/R)-\arccos(1/(1+\eps))$. 两端极限给 $\pi/2$. 收敛性由1附近的 $O((x-1)^{-1/2})$ 与无穷远的 $O(x^{-2})$ 保证.
\end{solution}
这个例子保留原课两个分支点的分支选择与收敛核查, 采用显式原函数完成数值计算. 钥匙孔轮廓在上一节已完整使用; 此例的数值计算采用原函数, 不另宣称完成含0的割线闭路计算.

\originalsection{习题与解答}
\begin{exercise}
求 $\int_\R\cos(tx)/(x^2+a^2)^2\dd x$, $a>0$.
\end{exercise}
\begin{solution}
对前面 $\pi\ee^{-a|t|}/a$ 关于 $a$ 求导. 在任意 $a$ 的小闭邻域, 导数被一个可积有理函数控制, 所以积分与导数可交换. 因 $\partial_a(x^2+a^2)^{-1}=-2a(x^2+a^2)^{-2}$, 结果为 $\pi(1+a|t|)\ee^{-a|t|}/(2a^3)$.
\end{solution}
\begin{exercise}
利用钥匙孔积分求 $\int_0^\infty x^{\alpha-1}/(1+x)^2\dd x$, $0<\alpha<1$.
\end{exercise}
\begin{solution}
分部积分对 $x^\alpha/(1+x)$ 求导, 两端值均为0, 得 $\int x^\alpha/(1+x)^2=\alpha I_\alpha$. 又 $I_\alpha=\int x^{\alpha-1}/(1+x)^2+\int x^\alpha/(1+x)^2$, 因而结果为 $(1-\alpha)\pi/\sin(\pi\alpha)$.
\end{solution}

\originalchapter{辐角原理与零点计数}
\sourceinfo{18.04 Topic 11 全部正文; 18.112 L12--L14. Nyquist 方向约定与稳定性区别为编者补充}
\originalsection{绕数}
\begin{definition}
闭路径 $\gamma$ 不经过 $a$ 时, 定义绕数 $\Ind(\gamma,a)=(2\pi\ii)^{-1}\int_\gamma\dd z/(z-a)$.
\end{definition}
\begin{proposition}\label{prop:index}
绕数是整数, 等于沿路径连续辐角的总增量除以 $2\pi$. 它在 $\C\setminus\gamma$ 每个连通分支上恒定, 在无界分支为0, 对不经过 $a$ 的同伦不变.
\end{proposition}
\begin{proof}
把参数区间分成有限小段, 使每段 $\gamma-a$ 的像在有全纯对数的小圆盘内. 在相邻段用 $2\pi\ii$ 的整数倍调整对数, 使其首尾衔接. 这样得到连续对数值 $\log|\gamma(t)-a|+\ii\theta(t)$. 逐段积分 $1/(z-a)$, 实部首尾相消, 虚部为 $\ii(\theta(b)-\theta(a))$. 闭路两端指数相同, 增量是 $2\pi$ 的整数倍.

积分核对 $a$ 连续且曲线距离非零, 所以绕数对 $a$ 连续; 连续整数值函数在连通分支恒定. 对足够大 $|a|$, 积分估计趋零, 因为整数只能最终为零, 无界分支全为零. 同伦不变性来自第3章应用于穿孔平面上的 $1/(z-a)$.
\end{proof}
标准逆时针简单闭曲线内部绕数1, 外部0. 例如单位圆绕两遍, 曲线的几何像仍是一个圆, 绕数却是2. 计数依赖路径的参数与方向, 不只是图上画出的集合.

\originalsection{辐角原理}
亚纯函数指除孤立极点外全纯的函数. 在曲线及其内部邻域亚纯且曲线上没有零点或极点时, 内部只有有限个零点和极点: 否则紧性给域内聚点, 违背零点或极点的孤立性.
\begin{theorem}[辐角原理]\label{thm:argument}
设 $f$ 在正向简单闭曲线 $C$ 及其内部的开邻域亚纯, 不恒为零, 且在 $C$ 无零、无极点. 设内部零点总阶数为 $N$, 极点总阶数为 $P$, 则
\[
\frac1{2\pi\ii}\int_C\frac{f'(z)}{f(z)}\dd z=N-P=\Ind(f\circ C,0).
\]
\end{theorem}
\begin{proof}
在零点 $a$ 写 $f=(z-a)^mg$, $g(a)\ne0$, 则 $f'/f=m/(z-a)+g'/g$, 留数为 $m$. 在 $m$ 阶极点同理留数为 $-m$. 其他点全纯, 留数定理给第一式. 换元沿像路径的积分恰为 $\int_Cf'/f$, 所以等于像路径绕数. 图像可能自交, 但绕数定义仍适用.
\end{proof}
\begin{theorem}[Rouché]\label{thm:rouche}
若 $f,g$ 在简单闭曲线及其内部邻域全纯, 且在边界上 $|g|<|f|$, 则 $f$ 与 $f+g$ 在内部的零点总数按重数相同.
\end{theorem}
\begin{proof}
对 $0\le t\le1$, $f+tg$ 在边界非零. 辐角原理的积分是连续于 $t$ 的整数函数: 分母在紧的边界参数和 $t$ 区间上有统一正下界. 所以零点计数不变.
\end{proof}
\begin{example}
求 $z^5+3z+1$ 在单位圆内的零点数, 以及在 $|z|<2$ 内的零点数.
\end{example}
\begin{solution}
单位圆上 $|z^5+1|\le2<|3z|=3$, 故与 $3z$ 一样, 有1个零点. 半径2圆上 $|3z+1|\le7<|z^5|=32$, 故有5个零点. 两条边界都无零点, 因此剩下4个零点位于 $1<|z|<2$.
\end{solution}

\originalsection{局部映射}
\begin{theorem}\label{thm:localmap}
非恒定全纯 $f$ 在 $a$ 附近若 $f-f(a)$ 有 $m$ 阶零点, 则对足够小的圆盘及足够接近 $f(a)$ 的每个 $w$, 方程 $f(z)=w$ 在圆盘内按重数有 $m$ 个解. 因而非常数全纯映射把开集映为开集. 若 $f'(a)\ne0$, 则在 $a$ 的某邻域一一且具有全纯逆映射.
\end{theorem}
\begin{proof}
选半径 $r$ 使 $a$ 是闭盘内 $f-f(a)$ 的唯一零点, 边界最小模 $\delta>0$. 对 $|w-f(a)|<\delta$, Rouché 比较 $f-f(a)$ 与常数 $f(a)-w$, 得 $m$ 个根, 所以像包含中心为 $f(a)$ 的小圆盘.

若 $m=1$, 对足够小闭盘, $f'$ 无零点. 上述小像圆盘中的每个 $w$ 都有唯一且简单的根; 原域中取该小圆盘的原像分支, 得一一对应. 逆映射连续性可由紧性验证: 任意 $w_n\to w$, 相应根的子序列聚点必须是唯一根, 所以整个序列收敛. 令 $z'=f^{-1}(w+h)$, 则
\[
\frac{f^{-1}(w+h)-f^{-1}(w)}h
 =\left(\frac{f(z')-f(z)}{z'-z}\right)^{-1}\longrightarrow\frac1{f'(z)}.
\]
故逆映射全纯.
\end{proof}
这里 ``按重数有 $m$ 个'' 不保证 $m$ 个不同点. 若 $m>1$, 更可写 $f-f(a)=h^m$, 其中 $h=(z-a)g^{1/m}$, 在小盘对非零 $g$ 选择对数后取根, 得 $h'(a)\ne0$. 因而局部行为就是一个保角坐标后的 $m$ 次幂映射. 非零全纯函数在单连通域存在对数的证明将在第13章给出; 小盘情形也可用对 $g(a)$ 附近的幂级数分支直接构造.

\originalsection{Nyquist 判据}
对于有理传递函数 $G(s)$, 非零负反馈系数 $k$ 的闭环函数为 $G_{\rm cl}=G/(1+kG)$. 令 $F=1+kG$. 为清楚固定方向, 取右半平面的正向大半圆边界 $C_R$: 虚轴从 $+\ii R$ 向 $-\ii R$ 下行, 再沿右半圆从底到顶. 区域始终在左侧. 若 $G$ 无虚轴极点且 $F$ 无虚轴零点, 半径足够大时辐角原理给
\[
Z=P+\Ind(1+kG(C_R),0)=P+\Ind(kG(C_R),-1),
\]
其中 $Z$ 为右半平面 $F$ 的零点数, $P$ 为其极点数; $k\ne0$ 时一般等于 $G$ 的右半平面极点数. 有理函数写成既约 $G=N/D$ 后, 闭环分母为 $D+kN$, 与分子 $N$ 无公共根, 因此零点计数确实对应闭环极点. 若有隐藏状态的精确消去, 输入输出极点计数不能自动替代内部稳定性检查.

当 $G$ 严格真有理时, 大半圆上的 $G\to0$, 像曲线该段收缩到0, 所以主要图形来自虚轴频率响应. 本方向下 $\omega$ 从正向负变化; 若惯常图采用从负向正的参数, 绕数要换符号. 只说 ``绕 $-1$ 几圈'' 而不声明方向, 判据没有确定数值.
\begin{example}
设 $G(s)=1/(s-1)$, $k>0$. 判断闭环稳定性.
\end{example}
\begin{solution}
闭环为 $1/(s+k-1)$, 极点 $1-k$. 所以 $k>1$ 时严格位于左半平面, $0<k<1$ 时有1个右半平面极点, $k=1$ 时位于虚轴, 不满足本判据的无边界零点条件. 从绕数看, 开环有 $P=1$, 稳定时要求 $\Ind(kG(C_R),-1)=-1$. 这与直接代数计算一致.
\end{solution}
对严格真有理且既约的因果系统, 所有极点严格在左半平面保证冲激响应可积, 因而有界输入产生有界输出. 虚轴上的简单极点只给某些自由解有界, 并不保证此种输入输出稳定性. 延迟使传递函数非有理, 有限大曲线辐角原理仍成立, 但还需核查无穷远行为及无穷多个极点的计数范围, 第12章给具体例子.

\originalsection{一般 Cauchy 与加权留数}
\begin{definition}
闭路径 $\gamma\subset\Omega$ 若对每个 $a\notin\Omega$ 都有 $\Ind(\gamma,a)=0$, 称相对于 $\Omega$ 零同调. 这只涉及本书已定义的整数绕数, 无需先修代数拓扑.
\end{definition}
\begin{theorem}\label{thm:generalcauchy}
若 $f$ 在开集 $\Omega$ 全纯, $\gamma$ 为其中零同调的闭路径, 则 $\int_\gamma f\dd z=0$; 对 $w\in\Omega\setminus\gamma$, 有
\[
\frac1{2\pi\ii}\int_\gamma\frac{f(z)}{z-w}\dd z=\Ind(\gamma,w)f(w).
\]
\end{theorem}
\begin{proof}
采用18.112 L13的 Dixon 路线并补足正则性. 定义对称差商 $g(w,z)=(f(z)-f(w))/(z-w)$, 对角上补 $f'(w)$. 局部 Taylor 展开或沿小圆盘的线段公式 $g(w,z)=\int_0^1f'(w+t(z-w))\dd t$ 说明联合连续; 固定一个变量时它全纯, 对角可去.

令 $\Omega_0=\{w\notin\gamma:\Ind(\gamma,w)=0\}$, 是开集. 在 $\Omega$ 定义 $H(w)=(2\pi\ii)^{-1}\int_\gamma g(w,z)\dd z$, 在 $\Omega_0$ 定义 $H(w)=(2\pi\ii)^{-1}\int_\gamma f(z)/(z-w)\dd z$. 交集上两式相同, 因差为 $f(w)\Ind(\gamma,w)=0$. 零同调保证 $\Omega\cup\Omega_0=\C$, 所以得到全平面函数.

在 $\Omega_0$ 核与路径有正距离, 可以微分积分, 得全纯. 在 $\Omega$ 对小三角形积分, 联合连续使有限路径的双积分可换序, 内层 $g$ 的三角形积分为零; Morera 给全纯. 所以 $H$ 整. 大 $|w|$ 时使用 $\Omega_0$ 的定义, 积分估计给 $H(w)=O(1/|w|)\to0$. Liouville 使 $H\equiv0$, 得加权公式. 再对全纯函数 $F(z)=(z-w_0)f(z)$ 应用该公式, $w_0\in\Omega\setminus\gamma$, 右边为零, 左边为 $(2\pi\ii)^{-1}\int_\gamma f$, 得零积分结论.
\end{proof}
若 $f$ 在 $\Omega$ 除有限个孤立奇点 $a_j$ 外全纯, 曲线不经过它们且相对 $\Omega$ 零同调, 则
\[
\int_\gamma f\dd z=2\pi\ii\sum_j\Ind(\gamma,a_j)\Res(f,a_j).
\]
证明可从每个奇点减去其 Laurent 主部. 主部在 $\C\setminus\{a_j\}$ 全纯, 负幂级数在曲线紧集上一致收敛; 指数 $-m$, $m\ne1$ 的项有全局原函数, 积分为零, $-1$ 次项给绕数. 减去主部后函数可补为 $\Omega$ 上全纯, 一般 Cauchy 使剩余积分为零. 此公式为自交或多次绕行时明确的计数规则.

\originalsection{习题与解答}
\begin{exercise}
证明 $z^4+6z+2$ 在单位圆内恰有一个零点.
\end{exercise}
\begin{solution}
边界上 $|z^4+2|\le3<6=|6z|$, 与 $6z$ 比较即得一个, 重数已计入.
\end{solution}
\begin{exercise}
设 $f$ 在闭单位圆盘邻域全纯, $|f(z)-z|<1$ 于圆周. 证明任意 $|w|<1-\max_{|z|=1}|f(z)-z|$ 都在 $f(\D)$ 内.
\end{exercise}
\begin{solution}
记边界最大差为 $M<1$. 对所给 $w$, $|(f-z)-w|\le M+|w|<1=|z|$. Rouché 比较 $f-w$ 与 $z$, 得盘内恰一个零点, 即有原像.
\end{solution}
\begin{exercise}
解释为什么 $\int_C f'/f$ 一般不为零, 即使 $f$ 在 $C$ 及其内部全纯.
\end{exercise}
\begin{solution}
若 $f$ 内部有零点, $f'/f$ 在那里有简单极点, 留数等于零点阶数, 不满足 Cauchy 零积分的全域全纯假设. 例如 $f(z)=z$, 单位圆积分为 $2\pi\ii$.
\end{solution}

\originalchapter{调和函数}
\sourceinfo{18.04 Topic 5 全部正文, Topic 10的 Dirichlet 问题; 18.112 L16的共轭、环域均值与 Poisson 公式}
\originalsection{调和共轭}
\begin{definition}
实函数 $u$ 在开集上为 $C^2$ 且满足 $\Delta u=u_{xx}+u_{yy}=0$, 称为调和函数. 若 $u+\ii v$ 全纯, 称 $v$ 为 $u$ 的调和共轭.
\end{definition}
全纯函数的实部和虚部调和. 例如 $u_{xx}+u_{yy}=v_{yx}-v_{xy}=0$, 混合偏导可交换是因为全纯函数已经由 Cauchy 理论证为无限次光滑. 共轭有方向: 若 $v$ 是 $u$ 的共轭, 则 $-u$ 是 $v$ 的共轭, 因为 $-\ii(u+\ii v)=v-\ii u$ 全纯.
\begin{theorem}\label{thm:conjugate}
单连通域上的任意调和函数 $u$ 都有全局调和共轭, 共轭只差实常数. 一般域上每点都有局部共轭, 因而调和函数自动无限次光滑.
\end{theorem}
\begin{proof}
令 $g=u_x-\ii u_y$. 由于 $u$ 为 $C^2$, 其各偏导连续, 且 $u_{xx}=-u_{yy}$、$u_{xy}=u_{yx}$ 恰为 $g$ 的 Cauchy--Riemann 方程. 所以 $g$ 全纯. 单连通时有原函数 $F=U+\ii V$. $F'=U_x-\ii U_y=g$ 给 $U_x=u_x,U_y=u_y$, 因而 $U-u$ 在域内为常数; 调整 $F$ 的实常数使 $U=u$. 两个共轭之差使全纯函数 $\ii(v_1-v_2)$ 纯虚, 因而恒定. 任意域局部可取小圆盘重复同一证明; $u$ 局部是全纯函数实部, 所以光滑.
\end{proof}
单连通不能删去. $u=\log|z|$ 在穿孔平面调和, 若有单值共轭, 则全纯 $F$ 的导数必须为 $1/z$, 与单位圆积分非零矛盾. 局部共轭就是辐角, 绕原点一周增加 $2\pi$.
\begin{example}
求 $u=x^2-y^2$ 和 $u=x/(x^2+y^2)$ 的共轭, 并核对等值线的正交性.
\end{example}
\begin{solution}
前者是 $z^2$ 的实部, 共轭为 $2xy$. 后者是 $1/z$ 的实部, 共轭为 $-y/(x^2+y^2)$. 一般情况下 $\nabla u\cdot\nabla v=u_xv_x+u_yv_y=0$. 当 $f'\ne0$ 时两个梯度都非零, 隐函数定理使等值线局部为光滑曲线, 两个法向量正交, 所以曲线正交. 在 $z^2$ 的0点梯度消失, $u=0$ 为两条对角线、$v=0$ 为两坐标轴, 都不是一条光滑曲线; 不能在那里宣称正交定理仍适用.
\end{solution}

\originalsection{均值与边值唯一性}
\begin{theorem}[调和最大值原理]\label{thm:harmonicmax}
圆盘内调和函数满足圆周平均值性质. 域内若有局部最大或局部最小, 则恒定. 有界域上连续到闭包的调和函数, 最大和最小均在边界取得. 因而连续 Dirichlet 边值问题至多有一个解.
\end{theorem}
\begin{proof}
在圆盘取共轭, 对全纯函数平均值取实部即可. 若 $u$ 在一点有局部最大值 $M$, 所有足够小圆上 $M-u\ge0$ 且平均为零, 连续性迫使圆上每点等于 $M$, 从而邻域恒定. 共轭函数实部在邻域恒定, 相应全纯函数导数为零, 恒等定理使其导数在域内恒为零; 或利用局部性质沿相交圆盘传播常数. 对 $-u$ 得最小值结论. 有界闭包紧, 极值存在, 非常数时只能在边界. 两个相同边值的解之差调和且边界为零, 最大最小值都为零.
\end{proof}
Dirichlet 问题是给定边界函数, 在内部寻找调和函数. 唯一性不自动给存在性; 下面在圆盘构造解. 若边界有跳跃, 也不能要求一个连续到整个闭包的函数同时满足两侧不同极限.

环域中圆周平均不一定恒定. 记 $A(r)=(2\pi)^{-1}\int_0^{2\pi}u(a+r\ee^{\ii t})\dd t$. 极坐标 Laplace 算子为 $\partial_r^2+r^{-1}\partial_r+r^{-2}\partial_t^2$. 对 $t$ 积分消掉周期二阶导数, 得 $(rA')'=0$, 因而 $A(r)=\alpha\log r+\beta$. 圆盘含中心时有限性迫使 $\alpha=0$, 恢复均值性质.

\originalsection{Poisson 公式}
\begin{theorem}[圆盘 Poisson 公式]\label{thm:poisson}
若 $u$ 在单位圆盘调和并连续到闭盘, 则对 $0\le r<1$,
\[
u(r\ee^{\ii\phi})=\frac1{2\pi}\int_0^{2\pi}
\frac{1-r^2}{1-2r\cos(\phi-t)+r^2}\,u(\ee^{\ii t})\dd t.
\]
反过来, 对连续圆周数据 $U$, 右式给调和函数且连续取得给定边值.
\end{theorem}
\begin{proof}
记 $z=r\ee^{\ii\phi}$. 核 $P_z(t)=(1-|z|^2)/|\ee^{\ii t}-z|^2$ 等于 $\Re[(\ee^{\ii t}+z)/(\ee^{\ii t}-z)]$. 所以固定有界可积 $U$, 积分是全纯函数实部, 在内部调和; 紧子盘内核与所有所需导数一致有界, 积分与微分可交换. 几何级数给 $P_z(t)=1+2\Re\sum_{n\ge1}z^n\ee^{-\ii nt}$, 积分可知核总质量为 $2\pi$, 且核非负.

当 $z$ 从盘内任意趋近圆周点 $\zeta_0$, 在 $|\ee^{\ii t}-\zeta_0|\ge\delta$ 的弧上, 分母最终有正下界而 $1-|z|^2\to0$, 核质量趋零. 近弧上利用 $U$ 在 $\zeta_0$ 的连续性, $|U-U(\zeta_0)|$ 小; 全核质量恒为 $2\pi$ 控制该误差. 因而边界极限为 $U(\zeta_0)$. 对连续 $U$, 每点都成立并给连续到闭盘的解. 原先的 $u$ 与该构造具有同一边值, 最大值唯一性给相等. 本证明不需要先假设 $u$ 在边界外全纯, 避免把连续边界条件误加强.
\end{proof}
若 $U$ 只有有限个跳点且分段连续, 同样在每个连续边界点收敛到 $U$. 在跳点沿径向趋近时, 偶对称的核给左右极限平均; 任意切向路径可能有不同极限, 因而不能将径向结论误写为任意方向结论.

上半平面的对应核为 $u(x+\ii y)=\pi^{-1}\int_\R yU(t)/((x-t)^2+y^2)\dd t$, $y>0$. 对有界连续数据积分存在. 对内部紧集, 实核及其二阶偏导在大 $|t|$ 时统一为 $O((1+t^2)^{-1})$, 乘有界 $U$ 后可积, 所以微分可移入积分; 每个实核对 $x,y$ 调和, 因而积分调和. 核总质量1, 且 $y\downarrow0$ 时远离 $x$ 的质量趋零, 同一论证给边值. 不加有界性或无穷远控制时, 上半平面的边界零值不保证唯一: $u=y$ 就是一个非零解.
\begin{example}
在单位圆盘求边值 $U(\ee^{\ii t})=3+2\cos2t-\sin3t$ 的调和延伸.
\end{example}
\begin{solution}
Poisson 核的 Fourier 展开使第 $n$ 次频率乘以 $r^n$, 所以 $u=3+2r^2\cos2\phi-r^3\sin3\phi$. 也可写 $u=\Re(3+2z^2+\ii z^3)$, 直接核对边值并用唯一性.
\end{solution}
\begin{example}
求上半平面的有界调和函数, 边值在 $x<0$ 为1、$x>0$ 为0.
\end{example}
\begin{solution}
取 $0<\arg z<\pi$, $u(z)=\arg z/\pi=\Re[\Log z/(\pi\ii)]$. 在左右两半轴分别趋于1、0, 且 $0<u<1$. 0处是数据跳点, 沿垂线趋近为 $1/2$, 沿射线角度 $\theta$ 趋近为 $\theta/\pi$. 多个阶梯边值也可叠加这些角函数: 断点 $a_j$ 处左值减右值为 $d_j$, 则 $u=c_{\rm right}+\sum_j d_j\arg(z-a_j)/\pi$.
\end{solution}

\originalsection{习题与解答}
\begin{exercise}
求环域 $1<|z|<R$ 内调和函数, 内边界值0、外边界值1, $R>1$.
\end{exercise}
\begin{solution}
$u=\log|z|/\log R$ 调和并满足边值. 环域有界, 最大值原理给唯一性. 它一般没有全局单值共轭, 因为共轭应为 $\arg z/\log R$.
\end{solution}
\begin{exercise}
若 $u$ 是平面上有上界的调和函数, 证明它恒定.
\end{exercise}
\begin{solution}
平面单连通, 存在整 $f$ 使 $u=\Re f$. 则 $\ee^f$ 有界整, Liouville 给恒定, 指数无零点所以 $f'=0$.
\end{solution}
\begin{exercise}
证明全纯 $T$ 与调和 $u$ 的复合 $u\circ T$ 仍调和, 并写出 Laplace 变换公式.
\end{exercise}
\begin{solution}
局部写 $u=\Re f$, 则 $u\circ T=\Re(f\circ T)$ 调和. 直接二元链式法则与 $T$ 的方程也给 $\Delta(u\circ T)=|T'|^2(\Delta u)\circ T$. 当 $T'\ne0$ 且有逆时, 这使调和边值问题可在两域间双向转移.
\end{solution}

\originalchapter{复势与平面流动}
\sourceinfo{18.04 Topic 6 全部正文; Topic 10的圆定理. 多连通域的周期障碍及流线计算为编者补充}
\originalsection{模型与复速度}
设平面定常速度场为 $V=(p(x,y),q(x,y))$, $p,q$ 为 $C^1$. 定常指不随时间改变; 不可压且密度恒定、无局部源汇给 $p_x+q_y=0$; 无旋给 $q_x-p_y=0$. 后两者分别是散度与旋度条件. 它们是所选模型的假设, 不是所有真实流动都满足的普遍规律.

散度定理的二维形式解释第一式: 正向闭边界上净向外通量为 $\int(p\dd y-q\dd x)=\iint(p_x+q_y)\dd x\dd y$. 无旋条件同理由环流 Green 公式说明. 穿孔域上的涡流可在每一点旋度为零, 却有围孔环流; 原点不在模型域内, 所以与局部无旋不矛盾.
\begin{definition}
全纯函数 $\Phi=\phi+\ii\psi$ 若满足 $V=\nabla\phi$, 称为速度场的复势. $\phi$ 是速度势, $\psi$ 是流函数. 复速度定义为 $\Phi'=p-\ii q$, 注意虚部的负号.
\end{definition}
\begin{theorem}
全纯复势产生不可压、无旋速度场. 反之, 在单连通域上, 任意 $C^1$ 且不可压、无旋的速度场都有复势.
\end{theorem}
\begin{proof}
全纯时 $\phi$ 调和, 所以 $\operatorname{div}\nabla\phi=0$, 混合偏导相等给旋度0, 方程又给 $\Phi'=\phi_x-\ii\phi_y$. 反向令 $g=p-\ii q$. 两条流动方程恰好是 $g$ 的 Cauchy--Riemann 方程, 所以 $g$ 全纯. 在单连通域取原函数 $\Phi$, 由 $\Phi'=\phi_x-\ii\phi_y=p-\ii q$ 得所需速度场.
\end{proof}
流体轨迹 $z(t)=(x(t),y(t))$ 满足 $(x',y')=V$. 沿轨迹 $\dd\psi/\dd t=\nabla\psi\cdot\nabla\phi=0$, 因而轨迹位于 $\psi$ 等值线上. 若 $V\ne0$, 等值线光滑且正是流线. 停滞点满足 $\Phi'=0$; 此处等值集可能交叉, 单纯看等值图不能确定流线的平滑延伸.

\originalsection{源、涡与偶极}
均匀速度 $V=(U,0)$ 的势为 $\Phi=Uz$, $U$ 为实常数. 流线是水平线, 方向由 $U$ 的符号决定. $\Phi=z^2$ 给 $V=(2x,-2y)$, 流线为 $2xy=\text{常数}$, 原点为停滞点. 在实轴向两侧流出, 在虚轴向原点流入.

源势 $\Phi=c\Log z$, $c\in\R$, 给 $V=(cx/r^2,cy/r^2)$, $\phi=c\log r$, $\psi=c\theta$. $c>0$ 为向外源, $c<0$ 为汇. 围原点的正向圆通量为 $2\pi c$, 环流0. 复势依赖分支, 速度 $c/z$ 在穿孔平面单值. 流函数绕孔增量为 $2\pi c$, 所以不存在单值全局复势.

涡势选择 $\Phi=-\ii\kappa\Log z$, 给 $\phi=\kappa\theta$, $\psi=-\kappa\log r$, $V=(-\kappa y/r^2,\kappa x/r^2)$. $\kappa>0$ 时逆时针, 环流 $2\pi\kappa$, 通量0. 原课程的 $\ii\Log z$ 对应 $\kappa=-1$, 所以是顺时针涡. 圆周围有整体环流并不表示穿孔域内每点有局部旋度.

更一般地, 若单值复速度 $g=p-\ii q$ 在某域全纯, 则
\[
\int_C g(z)\dd z=\int_C(p\dd x+q\dd y)+\ii\int_C(p\dd y-q\dd x).
\]
实部为环流、虚部为向外通量. 有单值原函数的充要条件是所有这些复周期为零. 这准确给出单连通假设在复势存在性中的作用.

\begin{example}
分析 $\Phi=\Log(z-1)+\Log(z+1)$ 的速度和停滞点, 并说明远处流动.
\end{example}
\begin{solution}
这是位于 $\pm1$ 的两等强源, 复速度 $\Phi'=1/(z-1)+1/(z+1)=2z/(z^2-1)$. 仅在0处停滞. 远处 $\Phi'=2/z+O(z^{-3})$, 相当于总强度2的源. 在实轴外段 $q=0$, 在虚轴上 $p=0$, 可由对称性或直接代入 $z=\ii y$ 得到; 因而坐标轴是相应流线的等值集. 局部可把两对数相加为 $\Log(z^2-1)$ 加一个常数, 但不能未经分支核对将此等式当全域恒等式.
\end{solution}
\begin{example}
将位置 $\pm h$ 的源汇对取极限, 得到偶极势.
\end{example}
\begin{solution}
取 $\Phi_h=(A/(2h))(\Log(z-h)-\Log(z+h))$. 对远离0的紧集, Taylor 展开或对参数积分给 $\Log(z-h)-\Log(z+h)=-2h/z+O(h^3)$, 一致成立. 因而 $\Phi_h\to-A/z$, 复速度趋 $A/z^2$. 此局部势称为偶极或双极势; 两源汇总通量为零, 但留下有方向的流场.
\end{solution}

\originalsection{圆形障碍}
\begin{theorem}[Milne--Thomson 圆定理]
设 $f$ 在闭圆盘 $|z|\le a$ 的开邻域全纯, 并在外部所需域上作为原流复势. 定义
\[
\Phi(z)=f(z)+\overline{f(a^2/\bar z)},\qquad |z|>a.
\]
第二项在圆外全纯, 不引入圆外的新奇点; 圆周上 $\Im\Phi=0$, 所以圆周为不穿透的流线边界.
\end{theorem}
\begin{proof}
若 $f(z)=\sum c_nz^n$, 且收敛半径大于 $a$, 则第二项等于 $\sum\bar c_na^{2n}z^{-n}$, 在 $|z|>a$ 的紧集上一致收敛, 故全纯. 因原函数在盘内邻域全纯, 反射点始终位于其正常域, 所以圆外无新增奇点. $|z|=a$ 时 $a^2/\bar z=z$, 于是 $\Phi=f(z)+\overline{f(z)}$ 为实数. 在非停滞点, $\psi$ 恒定使速度切向; 停滞点速度为零同样不穿透.
\end{proof}
\begin{example}
求半径 $a>0$ 的圆柱外匀速流及其停滞点, 其中 $U>0$.
\end{example}
\begin{solution}
取 $f=Uz$, 得 $\Phi=U(z+a^2/z)$, $\Phi'=U(1-a^2/z^2)$. 外部远处趋向均匀流, 停滞点为边界上的 $\pm a$. 极坐标中 $\psi=U(r-a^2/r)\sin\theta$, 圆周 $r=a$ 是 $\psi=0$. 速度分量为 $V_r=U(1-a^2/r^2)\cos\theta$, $V_\theta=-U(1+a^2/r^2)\sin\theta$, 显示边界法向分量为零.
\end{solution}
\begin{center}
\begin{tikzpicture}[>=Stealth,scale=.85]
\draw[->,gray] (-4.2,0)--(4.4,0) node[right]{$x$};
\draw[->,gray] (0,-2.6)--(0,2.8) node[above]{$y$};
\draw[thick] (0,0) circle (1);
\foreach \b in {.35,.8,1.3}{
\draw[domain=20:160,samples=90,smooth,variable=\t]
plot ({((\b/sin(\t)+sqrt((\b/sin(\t))^2+4))/2)*cos(\t)}, {((\b/sin(\t)+sqrt((\b/sin(\t))^2+4))/2)*sin(\t)});
\draw[domain=20:160,samples=90,smooth,variable=\t]
plot ({((\b/sin(\t)+sqrt((\b/sin(\t))^2+4))/2)*cos(\t)}, {-((\b/sin(\t)+sqrt((\b/sin(\t))^2+4))/2)*sin(\t)});
}
\fill (-1,0) circle (1.5pt); \fill (1,0) circle (1.5pt);
\node[below left] at (-1,0){$-a$}; \node[below right] at (1,0){$a$};
\end{tikzpicture}
\end{center}
图取 $a=U=1$, 各条流线由 $\psi=\pm0.35,\pm0.8,\pm1.3$ 的方程绘出. 一般流线满足 $y(1-a^2/(x^2+y^2))=\psi/U$, 图形不替代已经核验的边界条件.

对源位于 $-2$ 的 $f=\Log(z+2)$, 单位圆定理给 $\Phi=\Log(z+2)+\overline{\Log(1/\bar z+2)}$. 可在盘内选择与实轴相容的分支, 第二项等于相应 $\Log(2+1/z)$. 其复速度为 $1/(z+2)-1/[z(2z+1)]$, 修正项在远处为 $O(z^{-2})$, 而原源项为 $O(z^{-1})$. 新奇点位于圆内, 外流中只保留原源.

\originalsection{匀速流叠加源}
\begin{example}
设 $U,Q>0$, 分析 $\Phi(z)=Uz+(Q/(2\pi))\Log z$ 的流与停滞点.
\end{example}
\begin{solution}
记 $c=Q/(2\pi)$, 复速度为 $U+c/z$, 速度场为 $(U+cx/r^2,cy/r^2)$. 近原点源项占优, 向外通量为 $Q$; 无穷远趋于向右匀速. 唯一停滞点为 $z=-c/U$. 流函数 $\psi=Uy+c\theta$, 在上半平面该停滞点极限值为 $c\pi$, 下半平面为 $-c\pi$. 这两条等值分支构成从停滞点向下游分开的分界流线, 与实轴上游的流线衔接; 对数分支变化只加常数, 不改变速度. 在实轴 $x<-c/U$ 速度向右, $-c/U<x<0$ 速度向左, 表示源与均匀流相抵的方向变化.
\end{solution}
位置 $\pm a$ 的有限源汇对可写 $c[\Log(z-a)-\Log(z+a)]$. 复速度 $2ac/(z^2-a^2)$, 无有限停滞点, 净通量为零; 远处为 $2ac/z^2+O(z^{-4})$. 取 $a\to0$ 并固定 $2ac$ 就得到本章偶极极限, 因而有限源汇与极限模型的关系已同时保留.

\originalsection{习题与解答}
\begin{exercise}
场 $V=(cy,0)$ 是否不可压、无旋? $c\ne0$.
\end{exercise}
\begin{solution}
散度0, 旋度为 $q_x-p_y=-c\ne0$, 是剪切流, 不能由全纯复势表示. 水平箭头不旋转成圆并不意味着旋度为零.
\end{solution}
\begin{exercise}
在圆柱外匀速势中加入 $-\ii\kappa\Log z$. 求边界停滞点存在条件.
\end{exercise}
\begin{solution}
径向速度仍为零, 边界切向速度为 $-2U\sin\theta+\kappa/a$. 当 $U>0$ 时, 存在边界停滞点当且仅当 $|\kappa|\le2Ua$, 方程为 $\sin\theta=\kappa/(2Ua)$. 严格不等时有两个, 等号时合并为一个.
\end{solution}

\originalchapter{保角映射}
\sourceinfo{18.04 Topic 10 全部正文; 18.112 L6--L7为主题背景. 交比、单射及边界方向验证为编者补充}
\originalsection{局部与全局}
全纯 $f$ 在 $a$ 满足 $f'(a)\ne0$ 时, 沿任意两条非零切向量 $v_1,v_2$, 像切向量为 $f'(a)v_1,f'(a)v_2$, 它们商不变, 所以有向夹角保持. 这称为局部保角. 若 $f$ 还一一映某域 onto 另一域, 称为两域的双全纯或保角同构. 逆函数全纯已由局部映射定理证明. 局部保角不保证全局单射, 如指数在整个平面.

非恒定全纯单射必有处处非零导数. 若 $f'(a)=0$, 局部映射的重数 $m\ge2$; 选接近但不等于 $f(a)$ 且避开临界值的 $w$, 得至少两不同原像. 更直接, 局部表达 $f-f(a)=h^m$, $h$ 双全纯, $m$ 次幂不单射.

\originalsection{Möbius 变换}
\begin{definition}
映射 $T(z)=(az+b)/(cz+d)$, $ad-bc\ne0$, 称为 Möbius 变换. 在扩充平面上补 $T(-d/c)=\infty$, $T(\infty)=a/c$ ($c\ne0$); $c=0$ 时无穷远映为无穷远.
\end{definition}
代数求逆给 $T^{-1}(w)=(dw-b)/(-cw+a)$, 所以扩充平面一一对应. 有限正常点 $T'=(ad-bc)/(cz+d)^2\ne0$; 在无穷远及极点使用倒数坐标同样为保角. 矩阵 $\begin{pmatrix}a&b\\c&d\end{pmatrix}$ 的非零数倍给同一映射, 复合对应矩阵相乘.

每个 Möbius 变换由平移、非零复数乘法与倒数 $1/z$ 合成. 若 $c\ne0$, 将式子写成 $a/c+(bc-ad)/(c(cz+d))$ 即得分解. 平移和乘法显然将圆、直线映为圆、直线; 倒数的情况可用统一方程证明.
\begin{proposition}
广义圆是扩充平面中的普通圆或带无穷远点的直线. Möbius 变换将广义圆映为广义圆.
\end{proposition}
\begin{proof}
广义圆可写 $A|z|^2+Bz+\bar B\bar z+C=0$, $A,C\in\R$, $|B|^2-AC>0$. $A=0$ 时是直线, $A\ne0$ 时配方得正半径圆. 代入 $z=1/w$ 并乘 $|w|^2$, 形式变为 $C|w|^2+\bar Bw+B\bar w+A=0$, 仍为同类方程. 其他两种基本变换同理保留此形式, 因而任意复合也保留.
\end{proof}

\begin{definition}
不同的四点的交比取约定 $[z,z_1;z_2,z_3]=((z-z_2)(z_1-z_3))/((z-z_3)(z_1-z_2))$, 无穷远项按极限解释.
\end{definition}
恒等式 $T(z)-T(w)=(ad-bc)(z-w)/((cz+d)(cw+d))$ 表明交比不变. 给定三不同源点及三不同像点, 存在唯一 Möbius 变换实现对应: 先用上述交比把源三点 $z_1,z_2,z_3$ 映为 $1,0,\infty$, 再用目标相应变换的逆. 唯一性可由固定 $0,1,\infty$ 的变换只能为恒等证明.

四不同点在同一广义圆上当且仅当交比为实数. 把其中三点所在广义圆映为实轴, 第四点的交比就是其新坐标, 于是结论成立. 这是以三点确定广义圆的解析表达.
\begin{example}
证明 $T(z)=(z-\ii)/(z+\ii)$ 将上半平面一一映为单位圆盘.
\end{example}
\begin{solution}
对 $z=x+\ii y$, $|z+\ii|^2-|z-\ii|^2=4y$, 所以 $y>0$ 恰对应 $|T(z)|<1$. 逆变换 $z=\ii(1+w)/(1-w)$, 对 $|w|<1$ 有 $\Im z=(1-|w|^2)/|1-w|^2>0$. 因而像和逆像都明确, 并非只看实轴映为圆周便断言内部方向.
\end{solution}

\originalsection{圆的反射与同心化}
实轴反射为 $z\mapsto\bar z$, 单位圆反射为 $z\mapsto1/\bar z$. 一般圆 $|z-a|=r$ 的反射为 $a+r^2/(\bar z-\bar a)$, 中心与无穷远互相对应. 其推导是先平移伸缩到单位圆, 反射, 再变换回来. 点与反射点在同一径向线上且到中心的距离乘积为 $r^2$.

对称点的内在描述是: 所有经过两点的广义圆都与给定广义圆正交. 对实轴, 经过共轭两点的圆心位于实轴, 因而正交. Möbius 变换保持广义圆和角, 所以把这种描述及对称点关系一起保持. 这也是反射的唯一性依据.
\begin{proposition}
两条互不相交的广义圆可由 Möbius 变换映为同心圆.
\end{proposition}
\begin{proof}
先取第一条普通圆上的一点映为无穷远, 使它成为直线. 第二条与该点不相交, 所以映为普通圆, 且与该直线不相交. 再用相似变换把圆变为单位圆, 直线变为 $\Re z=c$, 其中可取 $c>1$. 实数 $u=c+\sqrt{c^2-1}$ 和 $v=c-\sqrt{c^2-1}$ 满足 $uv=1$, $u+v=2c$, 因而同时在单位圆和该直线互为反射. 拉回原域得到两点同时关于原两圆对称.

用 Möbius 变换把这两点映为 $0,\infty$. 两个像广义圆都关于 $0,\infty$ 有对称关系, 因而都必须是中心0的普通圆; 直线的反射不会将有限点0映为无穷远. 两圆于是同心. 如果原本其中一条已经是直线, 直接从相似归一化开始即可.
\end{proof}
这一结论也覆盖一圆包住另一圆且不相交的情形. ``不相交'' 排除了相切; 相切时共同对称点合并, 同心化构造退化.

\originalsection{典型区域}
设 $h>0$, $0<\alpha\le2\pi$. 条带 $0<\Im z<h$ 经 $\ee^{\pi z/h}$ 一一映为上半平面, 再接 $T$ 可映为圆盘. 扇形 $0<\arg z<\alpha$ 经 $z^{\pi/\alpha}$ 映为上半平面, 必须使用辐角在 $(0,\alpha)$ 的分支. 一个多值幂公式, 只有指定域与分支后才是一条映射.
\begin{example}
把上半单位圆盘 $|z|<1,\Im z>0$ 一一映为上半平面.
\end{example}
\begin{solution}
$q=(1+z)/(1-z)$ 满足 $\Re q=(1-|z|^2)/|1-z|^2>0$, $\Im q=2\Im z/|1-z|^2>0$, 所以映为第一象限. 逆式 $(q-1)/(q+1)$ 也把第一象限映回该半盘. 再平方得到 $q^2$, 一一映为上半平面. 两端点 $\pm1$ 映为0和无穷远, 半圆弧与直径分别映为负、正实半轴.
\end{solution}
域 $\Re z<0,0<\Im z<\pi$ 经 $\ee^z$ 映为上半单位圆盘, 再用前例可映为上半平面. 这种组合方式把复杂区域的映射分成每一步可核验的一一变换.

设 $a>0$. Joukowski 变换 $J(z)=z+a^2/z$ 将 $|z|>a$ 映为 $\C\setminus[-2a,2a]$. 单射性来自 $J(z_1)=J(z_2)$ 蕴含 $(z_1-z_2)(1-a^2/(z_1z_2))=0$, 外域不允许 $z_1z_2=a^2$. 对任意不在割线的 $w$, 二次方程 $z^2-wz+a^2=0$ 的两根乘积 $a^2$, 恰有一根在圆外: 如果有根在圆周, 写 $z=a\ee^{\ii\theta}$ 则 $w=2a\cos\theta$ 在割线上; 剩余情形由乘积与连续性判定内外, 或利用倒数配对和上、下半平面分支确定. 在上半圆外域, $\Im J=y(1-a^2/|z|^2)>0$, 所以其像为上半平面. 圆周映为线段且上下两岸分别对应不同边界弧.

\originalsection{边值与流动的转移}
若 $T:\Omega\to\Omega'$ 双全纯, $u$ 为 $\Omega'$ 的调和函数, 则 $u\circ T$ 调和. 对可连续对应的边界弧, 边值一起拉回. 在单位圆盘内, 边值右半圆1、左半圆0的解可用 $w=\ii(1+z)/(1-z)$ 转到上半平面. $z=\pm\ii$ 映为 $\mp1$, 右半圆映为 $(-\infty,-1)\cup(1,\infty)$, 左半圆映为 $(-1,1)$. 因而
\[
u(z)=1-\frac{\arg(w-1)-\arg(w+1)}\pi,
\qquad w=\ii\frac{1+z}{1-z},\quad 0<\arg(w\pm1)<\pi.
\]
直接看边界各区间上的角差可核实数值; 中心 $z=0,w=\ii$ 给 $u=1/2$, 与对称性吻合.

复势同样可拉回: 若 $\Phi(w)$ 表示 $w$ 域的流, 则 $\Psi(z)=\Phi(T(z))$ 在 $z$ 域有流线 $\Im\Phi(T(z))=\text{常数}$. 复速度为 $\Psi'=\Phi'(T)T'$, 所以变换会改变速度尺度, 并非只把原速度向量平移过去. 例如以 $T=z^2$ 把第一象限流动变到上半平面问题, 轴边界和半平面实轴由同一变换关联.

\originalsection{习题与解答}
\begin{exercise}
求将 $0,1,\infty$ 分别映为 $1,\infty,0$ 的 Möbius 变换.
\end{exercise}
\begin{solution}
无穷远映为0要求 $a=0$, 1映为无穷要求 $d=-c$, 0映为1要求 $b=d=-c$. 所以 $T(z)=-1/(z-1)=1/(1-z)$, 行列式非零.
\end{solution}
\begin{exercise}
求点 $1+\ii$ 在圆 $|z-2|=3$ 的反射.
\end{exercise}
\begin{solution}
代入 $2+9/(\bar z-2)$, 得 $2+9/(-1-\ii)=-5/2+9\ii/2$. 两点到中心的距离分别为 $\sqrt2$ 与 $9/\sqrt2$, 乘积9.
\end{solution}
\begin{exercise}
求从 $0<\arg z<\pi/3$ 到单位圆盘的双全纯映射, 并核对每步单射.
\end{exercise}
\begin{solution}
$z^3$ 在此扇形辐角变为 $(0,\pi)$ 且模可取任意正值, 一一映为上半平面. 接 $T$ 得 $(z^3-\ii)/(z^3+\ii)$, 上半平面与圆盘的双射已经验证.
\end{solution}

\originalchapter{Laplace 变换}
\sourceinfo{18.04 Topic 12 全部正文及 Topic 9的 Fourier 应用. 反演、冲激和稳定性条件为编者补充}
\originalsection{定义与全纯性}
\begin{definition}
局部分段连续函数 $f:[0,\infty)\to\C$ 的单侧 Laplace 变换为 $F(s)=\int_0^\infty f(t)\ee^{-st}\dd t$, 在积分绝对收敛的 $s$ 上定义. 若存在 $C,c$ 使 $|f(t)|\le C\ee^{ct}$ 于充分大的 $t$, 称 $f$ 为指数阶.
\end{definition}
只考虑 $t\ge0$, 并在需要讨论信号延迟时把 $f$ 延为 $t<0$ 时零. 有有限个跳点的分段连续函数按实、虚部作广义积分; 绝对可积指 $\int|f|<\infty$. 本章的积分换序和极限交换都由明确的绝对值控制进行. 为简洁, 第7章的 $L^1$ 记号仍表示绝对可积, 不要求读者先完成一般测度论.
\begin{theorem}
指数阶 $c$ 的 $f$ 在 $\Re s>c$ 有绝对收敛的全纯变换, 且 $F^{(n)}(s)=\int_0^\infty(-t)^nf(t)\ee^{-st}\dd t$.
\end{theorem}
\begin{proof}
在该半平面内的紧集取 $\Re s\ge c+\delta$, 尾部被 $Ct^n\ee^{-\delta t}$ 控制, 积分有限. 近0处局部分段连续保证有界可积. 截断积分为全纯函数, 尾估计在紧集上一致趋零; 全纯局部一致极限定理给全纯性和导数极限. 差商也可用同一 $\delta$ 缩小后的指数控制直接验证, 所以导数积分公式成立.
\end{proof}
例如 $F(s)=1/(s-a)$ 对应 $f(t)=\ee^{at}$, 积分仅在 $\Re s>\Re a$ 收敛. 有理表达虽可在 $s\ne a$ 继续定义, 左侧积分不会因此在其余区域突然收敛; 这是下一章解析延拓要区分的两件事.

\originalsection{运算规则}
线性由积分线性给出. 分部积分在 $f$ 连续且局部分段 $C^1$, $f,f'$ 满足足够的指数阶条件时给 $\mathcal L(f')=sF-f(0+)$. 反复应用,
\[
\mathcal L(f^{(n)})=s^nF(s)-\sum_{j=0}^{n-1}s^{n-1-j}f^{(j)}(0+).
\]
端点条件不可省略. 若函数在内部有跳点, 普通分段导数的变换还需要相应跳跃项, 或把导数解释为含冲激的广义导数.

乘时间对应 $\mathcal L(t^nf)=(-1)^nF^{(n)}$. 乘指数对应 $\mathcal L(\ee^{at}f)=F(s-a)$, 收敛半平面同时平移. 因果延迟 $f(t-\tau)\mathbf1_{t\ge\tau}$ 的变换为 $\ee^{-\tau s}F(s)$, $\tau\ge0$, 由 $u=t-\tau$ 换元得到. 积分规则为 $\mathcal L(\int_0^tf(u)\dd u)=F(s)/s$, 在足够右的半平面用导数规则或绝对换序证明.

另一条规则是 $\mathcal L(f(t)/t)(s)=\int_s^{+\infty}F(\zeta)\dd\zeta$, 右端沿水平射线 $\zeta=s+r$ 积分. 条件是 $f(t)/t$ 在0附近绝对可积且尾部为指数阶. 对此条件, 绝对换序给
\[
\int_0^\infty F(s+r)\dd r
 =\int_0^\infty f(t)\ee^{-st}\left(\int_0^\infty\ee^{-rt}\dd r\right)\dd t
 =\int_0^\infty\frac{f(t)}t\ee^{-st}\dd t.
\]
\begin{theorem}[卷积规则]
定义因果卷积 $(f*g)(t)=\int_0^tf(u)g(t-u)\dd u$. 若 $f,g$ 局部分段连续且为指数阶, 则在足够右的半平面 $\mathcal L(f*g)=FG$.
\end{theorem}
\begin{proof}
令 $v=t-u$, 积分区域 $0\le u\le t$ 变成 $u,v\ge0$. 绝对值双重积分不超过 $\int_0^\infty|f(u)|\ee^{-\sigma u}\dd u\int_0^\infty|g(v)|\ee^{-\sigma v}\dd v<\infty$, 其中 $\sigma=\Re s$ 足够大. 对有限矩形先用普通二重积分换序, 再用此尾控制扩展到无穷区域, 得
\[
\int_0^\infty\int_0^\infty f(u)g(v)\ee^{-s(u+v)}\dd u\dd v=F(s)G(s).
\]
这证明了卷积的变换公式.
\end{proof}
\begin{center}
\begin{tabular}{lll}
\toprule
函数 & 变换 & 一个绝对收敛范围\\
\midrule
$1$ & $1/s$ & $\Re s>0$\\
$t^n$, $n\in\Z_{\ge0}$ & $n!/s^{n+1}$ & $\Re s>0$\\
$\ee^{at}$ & $1/(s-a)$ & $\Re s>\Re a$\\
$\cos bt$, $b\in\R$ & $s/(s^2+b^2)$ & $\Re s>0$\\
$\sin bt$, $b\in\R$ & $b/(s^2+b^2)$ & $\Re s>0$\\
$\cosh bt$, $b\in\R$ & $s/(s^2-b^2)$ & $\Re s>|b|$\\
$\sinh bt$, $b\in\R$ & $b/(s^2-b^2)$ & $\Re s>|b|$\\
\bottomrule
\end{tabular}
\end{center}
这些公式由指数积分与递推分部积分得到, 不只是查表结果. 如 $\mathcal L(t^n)=n\mathcal L(t^{n-1})/s$, 起点 $\mathcal L(1)=1/s$. 表中的双曲函数范围采用 $|b|$, 避免负参数时把指数增长方向写错.

\originalsection{微分方程与系统函数}
线性系统指输入线性组合产生输出同一线性组合; 时不变指延迟输入使输出同样延迟. 常系数微分方程在零初始状态下定义因果线性时不变系统. 非零初值产生额外自由响应, 不能直接把整个输入输出映射称为线性零状态算子.
\begin{example}
解 $y''+3y'+2y=1$, $y(0)=y'(0)=0$, $t\ge0$.
\end{example}
\begin{solution}
变换后 $(s^2+3s+2)Y=1/s$, 所以 $Y=1/[s(s+1)(s+2)]=1/(2s)-1/(s+1)+1/[2(s+2)]$. 反演得 $y=1/2-\ee^{-t}+\ee^{-2t}/2$. 代回初值和方程可逐项核对. 零状态系统函数为 $H=1/[(s+1)(s+2)]$, 输出变换是 $HF$.
\end{solution}
一般 $P(D)y=Q(D)f$, $D=\dd/\dd t$, 在双方静止初始条件下变换为 $P(s)Y=Q(s)F$, 系统函数为 $H=Q/P$. 需要先约去公共因子, 并区分传递函数中的可见模式与实际状态方程中的隐藏模式.

冲激 $\delta_\tau$ 不当作普通函数, 而是作用于测试函数 $\varphi$ 给 $\varphi(\tau)$ 的对象. 因果单位冲激取 $\tau=0$ 的约定, $\mathcal L(\delta_0)=1$, $\mathcal L(\delta_\tau)=\ee^{-\tau s}$. 这一约定免除把端点冲激误算为一半的歧义. 对方程 $y''+3y'+2y=f$ 的因果 Green 函数为 $h(t)=(\ee^{-t}-\ee^{-2t})\mathbf1_{t\ge0}$. $h(0+)=0$, $h'(0+)=1$, 分布二阶导数产生单位冲激, 所以 $(D^2+3D+2)h=\delta_0$. 对一般输入 $y=h*f$, 卷积给零状态响应.

Fourier 方法对整条实轴上的平移不变方程同样有效. 例如衰减解 $-y''+a^2y=f$, $a>0$, 的变换为 $\widehat y=\widehat f/(\omega^2+a^2)$. 第7章算出的反变换核为 $K(x)=\ee^{-a|x|}/(2a)$, 因而 $y=K*f$. $K$ 连续, 导数从 $1/2$ 跳到 $-1/2$, 所以 $-K''+a^2K=\delta_0$. 对分段光滑、紧支撑 $f$, 可以直接微分两侧积分核验证方程和无穷远衰减. 这样原课的冲激、Green 函数与 Fourier 解微分方程形成同一条应用线索.

\originalsection{Fourier 乘子与因果解}
若 $f$ 连续、局部分段 $C^1$, $f,f'$ 绝对可积且 $f(\pm\infty)=0$, 在有限区间分部积分再取极限给 $\widehat{f'}(\omega)=\ii\omega\widehat f(\omega)$. 二阶时在相同的 $f'$ 条件下重复一次, 得 $\widehat{f''}=-\omega^2\widehat f$. 于是常系数微分算子 $P(D)$ 变成乘子 $P(\ii\omega)$. 有跳跃的函数则须把跳跃冲激计入导数, 不能直接套用普通导数规则.
\begin{example}
用 Fourier 方法求 $y''+8y'+7y=\ee^{-2t}\mathbf1_{t\ge0}$ 的因果零状态解.
\end{example}
\begin{solution}
外力的变换为 $1/(2+\ii\omega)$, 因而候选 $\widehat y=1/[(1+\ii\omega)(2+\ii\omega)(7+\ii\omega)]$. 该函数在实轴绝对可积, 反演时对 $t>0$ 在上半平面闭合, 对 $t<0$ 在下半平面闭合. 所有极点在 $\omega=\ii,2\ii,7\ii$, 下半平面没有极点; 圆弧积分由 $O(|\omega|^{-3})$ 与指数衰减趋零. 部分分式为
\[
\widehat y=\frac1{6(1+\ii\omega)}-\frac1{5(2+\ii\omega)}
 +\frac1{30(7+\ii\omega)}.
\]
每项的留数反演给相应因果指数, 所以 $y(t)=(\ee^{-t}/6-\ee^{-2t}/5+\ee^{-7t}/30)\mathbf1_{t\ge0}$. 在0右侧 $y=y'=0$, 与负半轴零解连续衔接且一阶导数连续; 因而二阶导数没有额外冲激. 直接代入得右侧外力. 这既验证 Fourier 乘子推导, 又明确因果条件如何由极点位置体现.
\end{solution}
对 $y''+y=\ee^{-2t}\mathbf1_{t\ge0}$, 因果 Green 核为 $\sin t\mathbf1_{t\ge0}$, 得 $y(t)=(\ee^{-2t}-\cos t+2\sin t)\mathbf1_{t\ge0}/5$. 它在正无穷不衰减, 不能假设普通绝对可积 Fourier 变换存在. 纯主值乘子 $\operatorname{PV}(1/(1-\omega^2))$ 对应对称 Green 核 $\sin|t|/2$, 会给另一个边界条件下的解. 原课实轴极点例的主值轮廓与因果选择在这里作此区分; 本书用因果卷积给完整解, 不把主值选择偷换成唯一的因果反演.

\originalsection{反演与唯一性}
对 $c$ 大于指数阶, 将 $g(t)=\ee^{-ct}f(t)$ 在负半轴补零, 则 $g\in L^1$ 且 $\widehat g(\omega)=F(c+\ii\omega)$. 因而第7章 Gaussian 反演在每个连续点 $t>0$ 给严格版本
\[
f(t)=\lim_{\eta\downarrow0}\frac1{2\pi\ii}
\int_{c-\ii\infty}^{c+\ii\infty}F(s)\ee^{st}\ee^{\eta(s-c)^2}\dd s.
\]
在这条竖线上, 正则化因子为 $\ee^{-\eta\omega^2}$. 若 $F(c+\ii\omega)$ 可积, 可去掉它, 得普通绝对收敛的 Bromwich 积分. 对适当分段 $C^1$ 函数, 也有对称截断的条件收敛版本; 对有理变换下面直接证明, 不把一般连续性当作足以保证所有普通 Fourier 积分收敛的条件.

若两个连续指数阶函数的 Laplace 变换相同, 其差的 Fourier 变换在这条竖线为零, Gaussian 反演给差在所有 $t>0$ 为零, 连续性再给 $t=0$ 也为零. 分段连续函数只保证连续点相同; 在有限个点任意改值不会改变积分, 因而不宣称对每个点无条件唯一.
\begin{theorem}[有理变换的留数反演]\label{thm:laplaceres}
若 $F$ 为严格真有理函数, 所有极点 $a_j$ 均在 $\Re s<c$, 则对 $t>0$,
\[
f(t)=\sum_j\Res\bigl(\ee^{st}F(s),a_j\bigr)
 =\lim_{R\to\infty}\frac1{2\pi\ii}\int_{c-\ii R}^{c+\ii R}\ee^{st}F(s)\dd s,
\]
且 $\mathcal L(f)=F$ 于足够右的半平面.
\end{theorem}
\begin{proof}
部分分式分解写 $F=\sum_{j,m}A_{jm}(s-a_j)^{-m}$. 分部积分或函数表给 $\mathcal L[t^{m-1}\ee^{a_jt}/(m-1)!]=(s-a_j)^{-m}$. 另一方面 $\ee^{st}$ 在 $a_j$ 的 Taylor 系数使该项留数正是 $A_{jm}t^{m-1}\ee^{a_jt}/(m-1)!$. 有限求和证明第一表达与变换关系.

对竖线积分, 用左闭合矩形, 顶底高度 $\pm R$, 左边 $\Re s=-R$. 因 $F$ 严格真有理, 在足够大 $|s|$ 有 $|F(s)|\le C/|s|$. 顶底边上的积分模至多为 $C\int_{-R}^c\ee^{tx}\dd x/R\le C\ee^{tc}/(tR)$, 左边至多 $2C\ee^{-tR}$. 三边趋零. 矩形正向, 右边由下向上, 留数定理使右边截断积分趋向上述留数和.
\end{proof}
该公式不能仅凭有限极点个数便用于任何增长的亚纯函数; 圆外的衰减或关闭轮廓的估计不可删去. 若变换有多项式部分, 其逆含冲激及其导数, 也不全是普通时间函数.
\begin{example}
求 $F(s)=1/(s^2+b^2)^2$ 的逆变换, $b>0$.
\end{example}
\begin{solution}
由 $\mathcal L(\sin bt)=b/(s^2+b^2)$ 对 $b$ 求导并组合, 或直接算两二阶极点, 得 $f(t)=(\sin bt-bt\cos bt)/(2b^3)$. 逐项变换为 $1/[2b^2(s^2+b^2)]-(s^2-b^2)/[2b^2(s^2+b^2)^2]$, 合并即原式.
\end{solution}

\originalsection{延迟反馈}
负反馈系统满足 $y=H*(f-k y)$, 变换后 $Y=HF-kHY$, 所以闭环 $H/(1+kH)$. 延迟 $\tau$ 的反馈将分母改成 $1+k\ee^{-\tau s}H(s)$. 即便 $H$ 有理, 新函数一般有无穷多个极点.

取 $H=1$, $\tau>0$, $k>0$. 极点由 $\ee^{-\tau s}=-1/k$ 给
\[
s_n=\frac{\log k}{\tau}+\frac{(2n+1)\pi\ii}{\tau},\qquad n\in\Z.
\]
符号可通过取模核对. $k<1$ 时全在左半平面, $k>1$ 时全在右半平面. 更强的输入输出结论可从时域直接证明: 对足够右的 $s$, 几何展开给 $Y=\sum_{n\ge0}(-k)^n\ee^{-n\tau s}F$, 故 $y(t)=\sum_{0\le n\le t/\tau}(-k)^nf(t-n\tau)$. 每个有限 $t$ 只有有限项.

若 $0<k<1$, 则 $|y(t)|\le\|f\|_\infty/(1-k)$, 得有界输入有界输出. 若 $k>1$, 可取持续时间小于 $\tau$ 的有界脉冲输入, 每次延迟复制产生模为 $k^n$ 的脉冲, 故不稳定. $k=1$ 时取不同延迟区间上符号交替的有界输入, 在相同相位点各项同向相加, 输出随项数增长, 也不能称为有界输入有界输出稳定. 这里没有靠 ``极点看起来都在左边'' 代替无限维系统所需的估计.

\originalsection{习题与解答}
\begin{exercise}
解 $y'+2y=f(t)$, $y(0)=y_0$, 并区分自由响应和零状态响应.
\end{exercise}
\begin{solution}
变换给 $(s+2)Y=F+y_0$. 所以 $y=y_0\ee^{-2t}+\int_0^t\ee^{-2(t-u)}f(u)\dd u$. 第一项为自由响应, 第二项为卷积零状态响应. 求导可核对初值与方程.
\end{solution}
\begin{exercise}
若 $f(t)=\sin t$ ($t\ge0$), 求延迟两单位后的变换, 并求 $t\sin t$ 的变换.
\end{exercise}
\begin{solution}
延迟信号必须写 $\sin(t-2)\mathbf1_{t\ge2}$, 变换为 $\ee^{-2s}/(s^2+1)$. 时间乘法给 $-\dd[(s^2+1)^{-1}]/\dd s=2s/(s^2+1)^2$.
\end{solution}

\originalchapter{解析延拓与 Gamma 函数}
\sourceinfo{18.04 Topic 13 全部正文; 18.112 L18的 Beta、反射与乘积路线. 乘积、倍增和实轴 Stirling 证明为编者补充}
\originalsection{延拓与分支}
\begin{definition}
若 $\Omega\subset\Omega'$ 为域, $f$ 在 $\Omega$ 全纯, 而 $F$ 在 $\Omega'$ 全纯且限制到 $\Omega$ 等于 $f$, 则 $F$ 为 $f$ 到 $\Omega'$ 的解析延拓. 允许孤立极点时称为亚纯延拓.
\end{definition}
恒等定理保证到一个固定连通域的延拓唯一, 却不保证延拓存在, 也不保证沿不同闭路的局部延拓最后给同一分支. 例如 $\Log z$ 沿绕0一周的路径逐盘延拓后增加 $2\pi\ii$. 主对数域与辐角 $(0,2\pi)$ 的对数域的交集分成上下半平面: 两分支在上半平面相同、在下半平面相差 $2\pi\ii$. 共同域不连通, 并不违背恒等定理.
\begin{proposition}\label{prop:logexist}
单连通域内无零点的全纯 $f$ 有全纯对数 $L$, 即 $\ee^L=f$, 因而也有任意给定正整数次根.
\end{proposition}
\begin{proof}
$f'/f$ 全纯, 取原函数 $H$. 则 $(f\ee^{-H})'=0$, 所以 $f\ee^{-H}=c\ne0$. 选一个数 $\ell$ 满足 $\ee^\ell=c$, 得 $L=H+\ell$. $\ee^{L/n}$ 为所需根. 两对数只差 $2\pi\ii$ 的整数倍, 选择初始值后唯一.
\end{proof}
若能沿域内任意路径解析延拓, 且沿同伦路径所得局部函数相同, 单连通性使其形成单值函数. 同伦路径相同的依据是用有限小圆盘覆盖同伦像, 在小网格上靠恒等定理逐格传播相同的局部函数. 这称为单值化原理; 其假设中的 ``沿每条路径可延拓'' 是实质条件, 不能由单连通自动保证任意函数可延到任何目标域.

\originalsection{Gamma 积分}
\begin{definition}
对 $\Re z>0$, 定义 $\Gamma(z)=\int_0^\infty t^{z-1}\ee^{-t}\dd t$, 其中 $t^{z-1}=\ee^{(z-1)\log t}$ 使用正实数对数.
\end{definition}
\begin{theorem}
积分在右半平面绝对收敛且全纯, 满足 $\Gamma(1)=1$, $\Gamma(z+1)=z\Gamma(z)$. 因而对整数 $n\ge0$, $\Gamma(n+1)=n!$. 它亚纯延拓到全平面, 恰在 $0,-1,-2,\ldots$ 有简单极点, 留数为 $(-1)^n/n!$.
\end{theorem}
\begin{proof}
在紧集上取 $0<\delta\le\Re z\le M$. 零点附近的绝对值由 $t^{\delta-1}$ 控制, 无穷远由 $t^{M-1}\ee^{-t}$ 控制, 任意阶 $z$ 导数增加的 $|\log t|^k$ 仍可积. 截断积分局部一致极限或受控差商给全纯性. 分部积分中 $t^z\ee^{-t}$ 在两端为零, 得递推式, 初值直接积分得到.

递推反写为 $\Gamma(z)=\Gamma(z+m+1)/[z(z+1)\cdots(z+m)]$. 对 $\Re z>-m-1$ 除分母零点外定义亚纯函数, 重叠域上由递推一致. 在 $z=-n$ 选 $m=n$, 分子趋 $\Gamma(1)=1$, 分母除 $z+n$ 外趋 $(-1)^nn!$, 所以简单极点留数 $(-1)^n/n!$. 非整数点选足够大 $m$ 即全纯, 因而没有其他极点.
\end{proof}
Gamma 延拓不是把发散的原积分赋一个普通积分值. 正实轴上 $\Gamma(x)>0$, 但其余复点的非零性仍需要证明, 下面的乘积会给出.

若 $\Re\alpha>0$, $\Re s>0$, 则 $\mathcal L(t^{\alpha-1})(s)=\Gamma(\alpha)s^{-\alpha}$, $s$ 的对数选右半平面分支. 先对实 $s>0$ 换元 $u=st$, 得式; 再对 $s$ 用恒等定理扩到右半平面. 在 $t=0$ 不必有普通函数有界性, $\Re\alpha>0$ 的局部可积性已足够. 原课较强的 $\Re\alpha>1$ 条件可放宽到这一最适范围. 例如 $\mathcal L(1/\sqrt{\pi t})=1/\sqrt s$.

\originalsection{Beta、反射与倍增}
定义 $B(z,w)=\int_0^1t^{z-1}(1-t)^{w-1}\dd t$, $\Re z,\Re w>0$. 绝对收敛使二重 Gamma 积分可换元 $x=r t,y=r(1-t)$, Jacobian 为 $r$, 得
\[
\Gamma(z)\Gamma(w)=\Gamma(z+w)B(z,w).
\]
这是恒等式形式, 即使暂未证明分母非零, 也没有贸然作除法. 当 $z,w$ 为正实数时分母正, 可写成商; 随后乘积非零性保证复参数商式同样有意义.

取实 $0<x<1$, $w=1-x$, 再换元 $t=u/(1+u)$, 得 $\Gamma(x)\Gamma(1-x)=\int_0^\infty u^{x-1}/(1+u)\dd u$. 第7章钥匙孔积分已给 $\pi/\sin\pi x$. 两边亚纯, 恒等定理从实区间推广, 所以
\[
\Gamma(z)\Gamma(1-z)=\frac\pi{\sin\pi z}
\]
作为亚纯恒等式成立, 在非整数点可直接取值. 负整数的极点与另一因子的有限值由递推精确协调.

倍增公式由 Beta 积分的对称性得到. 对 $\Re z>0$, 用 $t=(1+u)/2$ 并分两半,
\[
B(z,z)=2^{1-2z}\int_{-1}^1(1-u^2)^{z-1}\dd u
 =2^{1-2z}B(1/2,z).
\]
代入 Gamma 关系并消去正实参数下非零因子, 再解析延拓, 得
\[
\Gamma(z)\Gamma(z+1/2)=2^{1-2z}\sqrt\pi\,\Gamma(2z).
\]
$\Gamma(1/2)=\sqrt\pi$ 可先由 $t=u^2$ 得 $2\int_0^\infty\ee^{-u^2}\dd u$, 将 Gaussian 积分平方后以极坐标求得 $\pi$. 因而计算该常数不依赖尚待证明的倍增公式. 递推给 $\Gamma(3/2)=\sqrt\pi/2$.

\originalsection{Euler 乘积}
\begin{theorem}\label{thm:gammaproduct}
令 $\gamma=\lim_{N\to\infty}(\sum_{n=1}^N1/n-\log N)$, 则
\[
\frac1{\Gamma(z)}=z\ee^{\gamma z}\prod_{n=1}^\infty\left(1+\frac zn\right)\ee^{-z/n}.
\]
右侧在全平面局部一致收敛为整函数, 零点恰为 $0,-1,-2,\ldots$, 全为简单零点. 特别地 Gamma 在其有限正常点处从不为零.
\end{theorem}
\begin{proof}
先说明 Euler 常数存在. $H_N-\log N$ 由积分比较有下界, 且增量 $1/(N+1)-\log(1+1/N)<0$, 所以单调收敛.

对 $\Re z>0$, Beta 积分反复分部积分给 $B(z,N+1)=N!/[z(z+1)\cdots(z+N)]$. 令 $u=Nt$, 则
\[
N^zB(z,N+1)=\int_0^N u^{z-1}(1-u/N)^N\dd u\longrightarrow\Gamma(z).
\]
因为 $0\le(1-u/N)^N\le\ee^{-u}$, 在参数紧集上用 $u^{\delta-1}$ 与 $u^{M-1}$ 的分段控制, 极限局部一致. 倒数有限乘积为
\[
\frac{z(z+1)\cdots(z+N)}{N!N^z}
 =z\ee^{z(H_N-\log N)}\prod_{n=1}^N(1+z/n)\ee^{-z/n}.
\]
对固定紧集及足够大 $n$, 主对数展开给 $|\Log(1+z/n)-z/n|\le C_K/n^2$. 对数尾级数一致绝对收敛, 故乘积局部一致收敛; 远离列出的有限零因子时尾乘积非零. 因此极限整函数 $E$ 的零点恰为这些简单零点. 在右半平面, 上述有限 Gamma 近似与倒数乘积相乘恒为1, 两者都局部一致收敛, 得 $\Gamma E=1$, 同时证明 Gamma 无零点. 亚纯延拓使此恒等式遍及全平面, 给所述乘积.
\end{proof}
无穷乘积不能像有限乘积那样未经收敛论证便求导. 在避开零点的紧集, 局部一致极限定理保证有限乘积导数趋于极限导数, 分母又有正下界, 所以对数导数合法. 本乘积给
\[
\frac{\Gamma'(z)}{\Gamma(z)}=-\gamma-\frac1z+
\sum_{n=1}^\infty\left(\frac1n-\frac1{n+z}\right).
\]
尾项在紧集为 $O(n^{-2})$. 将 $z$ 与 $-z$ 的倒数乘积相乘, 再结合反射公式和递推, 得正弦乘积 $\sin\pi z=\pi z\prod_{n\ge1}(1-z^2/n^2)$. 在整数外取对数导数还得 $\pi\cot\pi z=1/z+\sum_{n\ge1}2z/(z^2-n^2)$, 与留数求和方法相互核对.

\originalsection{Stirling 的实轴版本}
\begin{theorem}
实数 $x\to+\infty$ 时, $\Gamma(x+1)\sim\sqrt{2\pi x}(x/\ee)^x$. 因而 $n!\sim\sqrt{2\pi n}(n/\ee)^n$.
\end{theorem}
\begin{proof}
从积分出发令 $t=xv$, 得
\[
\Gamma(x+1)=x^{x+1}\ee^{-x}\int_0^\infty
\ee^{-x h(v)}\dd v,\qquad h(v)=v-1-\log v.
\]
$h\ge0$, 仅在 $v=1$ 为零, 且 $h(1)=h'(1)=0,h''(1)=1$. 固定小 $\delta\in(0,1/2)$, 在 $|v-1|\le\delta$ 上有 $h(v)\ge c(v-1)^2$; Taylor 展开给 $xh(1+u/\sqrt x)\to u^2/2$. 换元后被 $\ee^{-cu^2}$ 控制, 将积分区间外补零, 控制收敛给局部积分乘 $\sqrt x$ 趋 $\int_\R\ee^{-u^2/2}\dd u=\sqrt{2\pi}$.

剩余两端不是只用点态极限忽略. $0<v\le1-\delta$ 时 $h(v)\ge h(1-\delta)>0$, 区间长度有限, 积分指数小. $v\ge1+\delta$ 时 $h'(v)=1-1/v\ge\delta/(1+\delta)=d>0$, 所以 $h(v)\ge h(1+\delta)+d(v-1-\delta)$, 尾积分至多 $\ee^{-xh(1+\delta)}/(xd)$. 两者乘 $\sqrt x$ 都趋零. 故总积分等价于 $\sqrt{2\pi/x}$, 代回即得.
\end{proof}
原课列出复参数渐近而略去证明. 复域版本需要固定对数分支及远离负实轴的扇区条件, 本书不据实轴证明声称其已成立; 对应拓展边界在来源表中保留.

\originalsection{习题与解答}
\begin{exercise}
计算 $\Gamma(5/2)$, $\Gamma(-1/2)$, 以及 Gamma 在 $-3$ 的留数.
\end{exercise}
\begin{solution}
递推给 $\Gamma(5/2)=3\sqrt\pi/4$, $\Gamma(-1/2)=-2\sqrt\pi$. 留数为 $(-1)^3/3!=-1/6$.
\end{solution}
\begin{exercise}
证明 $\int_0^1t^{x-1}(1-t)^n\dd t=n!/[x(x+1)\cdots(x+n)]$, $x>0$, $n\ge0$.
\end{exercise}
\begin{solution}
$n=0$ 时为 $1/x$. $n\ge1$ 时分部积分取 $\dd(t^x)/x$, 两端项为零, 得 $n/x$ 乘 $B(x+1,n)$. 递推即得式, 也说明 Euler 乘积近似不需调用未证明的 Beta 商式.
\end{solution}
\begin{exercise}
原积分 $\int_0^\infty\ee^{3t}\ee^{-zt}\dd t$ 能否在 $z=0$ 使用? 其解析延拓在0是多少?
\end{exercise}
\begin{solution}
原积分仅在 $\Re z>3$ 收敛, $z=0$ 发散. 在该半平面值为 $1/(z-3)$, 解析延拓到 $\C\setminus\{3\}$ 后在0值为 $-1/3$. 延拓值与发散积分不能混称.
\end{solution}

\originalchapter{正规族与 Riemann 映射}
\sourceinfo{拓展: 18.112 L12、L19、L20. L20仅为教材证明注记; 本章完整存在性证明与反射原理为编者补充}
\originalsection{Schwarz 引理}
\begin{theorem}[Schwarz]\label{thm:schwarz}
若 $f:\D\to\D$ 全纯且 $f(0)=0$, 则 $|f(z)|\le|z|$, $|f'(0)|\le1$. 任一非零点取等, 或导数模取等, 则 $f(z)=\ee^{\ii\theta}z$.
\end{theorem}
\begin{proof}
把 $g=f/z$ 在0补为 $f'(0)$, 得全纯函数. 在 $|z|=r<1$ 上 $|g(z)|\le1/r$, 最大模原理给内部同一界. 对固定 $z$ 令 $r\uparrow1$, 得 $|g(z)|\le1$. 若在某点模为1, 最大模原理使 $g$ 为单位模常数.
\end{proof}
圆盘自同构 $\phi_a(z)=(z-a)/(1-\bar az)$ 满足
\[
1-|\phi_a(z)|^2=\frac{(1-|a|^2)(1-|z|^2)}{|1-\bar az|^2}>0.
\]
其逆为 $(w+a)/(1+\bar aw)$. 任意圆盘自同构经 $\phi_a$ 归一化后固定0, 对函数和逆都用 Schwarz, 得它是旋转. 所以所有圆盘自同构形如 $\ee^{\ii\theta}\phi_a$.

对 $f:\D\to\D$ 全纯, 用 $\phi_{f(a)}\circ f\circ\phi_a^{-1}$ 应用 Schwarz, 得 Schwarz--Pick 估计
\[
\left|\frac{f(z)-f(a)}{1-\overline{f(a)}f(z)}\right|
\le\left|\frac{z-a}{1-\bar az}\right|,\qquad
\frac{|f'(a)|}{1-|f(a)|^2}\le\frac1{1-|a|^2}.
\]
定义长度元 $|\dd z|/(1-|z|^2)$, 对任意分段光滑路径应用第二式后积分, 得像路径长度不增加. 这给18.112中非 Euclid 圆盘解释的解析基础; 完整双曲几何公理及测地线理论不在本章范围.

\originalsection{正规族与 Hurwitz 定理}
\begin{definition}
域上全纯函数族称为正规, 若任意函数序列都有在每个紧集上一致收敛于有限全纯函数的子序列. 本书使用有限值版, 不把趋于无穷的亚纯正规族版本混入.
\end{definition}
\begin{theorem}[Montel]\label{thm:montel}
若全纯函数族在每个紧集上统一有界, 则正规.
\end{theorem}
\begin{proof}
固定紧集 $K\Subset\Omega$, 取 $d>0$ 使其闭 $2d$ 邻域仍为域内紧集, 函数族在其上统一有界 $M$. Cauchy 导数估计给闭 $d$ 邻域上 $|f'|\le M/d$. 对 $K$ 内距离小于 $d$ 的两点, 线段位于该邻域, 故 $|f(z)-f(w)|\le(M/d)|z-w|$. 全族于是等度连续.

域内取可数稠密点列, 对每点上的有界数列逐次取收敛子序列, 再取对角子序列, 得在所有稠密点收敛. 对任意紧集, 用有限个小圆盘覆盖, 每盘中心选附近稠密点, 等度连续使函数在任意点与该稠密点的差统一小. 稠密点有限, 所以对角子序列在那里最终互差小. 三角不等式使其在整个紧集一致 Cauchy, 复数完备性给一致极限. 选稠密点时使用 $K$ 的一个紧邻域, 不能错误要求稠密点必须属于 $K$ 本身. 第4章局部一致极限定理保证极限全纯.
\end{proof}
\begin{theorem}[Hurwitz]\label{thm:hurwitz}
域上无零点全纯 $f_n$ 局部一致收敛于 $f$, 则 $f$ 无零点或恒为零. 单射全纯函数的局部一致极限, 若非常数, 则仍单射.
\end{theorem}
\begin{proof}
若 $f$ 非恒零且在 $a$ 有零点, 选小圆使该点为盘内唯一零点, 边界 $|f|$ 有正下界. 对充分大 $n$, $|f_n-f|<|f|$ 于边界, Rouché 使 $f_n$ 有相同正数个零点, 矛盾.

对第二项, 固定 $w\in\Omega$, 考察 $g_n(z)=(f_n(z)-f_n(w))/(z-w)$, 在 $w$ 补为 $f_n'(w)$. 单射函数导数非零, 所以 $g_n$ 无零点. 导数的局部一致收敛与差商积分表达保证 $g_n$ 局部一致收敛于对应 $g$. 若 $f$ 非常数, $g$ 不恒为零, 故无零点. 因而 $f(z)\ne f(w)$ 对任何 $z\ne w$, 得单射.
\end{proof}

\originalsection{Riemann 映射定理}
\begin{theorem}[Riemann]\label{thm:riemann}
若 $\Omega\subsetneq\C$ 非空单连通域, 则存在双全纯映射 $f:\Omega\to\D$. 给定 $a\in\Omega$, 要求 $f(a)=0$ 和 $f'(a)>0$ 为正实数后, 映射唯一.
\end{theorem}
\begin{proof}
先构造到圆盘的一个单射全纯映射, 使极值族非空. 取 $b\notin\Omega$, 无零点函数 $z-b$ 在单连通域有平方根 $h$. 若 $h(z_1)=h(z_2)$, 平方得 $z_1=z_2$, 所以 $h$ 单射. 它的像不同时包含互为相反的两个点: 平方相等将迫使原像相同, 而 $h\ne0$. 取 $c=h(a)$, 开映射性给某个 $\eps>0$ 使 $D(c,\eps)\subset h(\Omega)$, 则 $D(-c,\eps)$ 与像不相交. 因而 $1/(h+c)$ 单射全纯且模不超过 $1/\eps$. 适当缩放使像严格在圆盘内, 再用圆盘自同构和旋转, 得归一化映射 $f_0(a)=0,f_0'(a)>0$.

令 $\mathcal F$ 为所有这样的归一化单射映射. 在 $a$ 的固定小圆上 Cauchy 估计给 $f'(a)\le M_0$, 所以导数上确界 $M$ 有限且 $M\ge f_0'(a)>0$. 取 $f_n'(a)\to M$, 族在域内由1统一有界, Montel 给局部一致极限 $f$. 导数收敛给 $f'(a)=M>0$, 因而非常数. 最大模原理使像严格在圆盘内, Hurwitz 使它单射, 所以 $f\in\mathcal F$ 并取得最大导数.

假若有 $b\in\D\setminus f(\Omega)$, 因 $f(a)=0$, $b\ne0$. $\phi_b\circ f$ 无零点, 在单连通域取全纯平方根 $h$; 它在圆盘内且单射. 记 $c=h(a)$, $|c|=\sqrt{|b|}$. 归一化 $g=\ee^{\ii\theta}\phi_c\circ h$, 旋转使 $g'(a)>0$. 链式法则给
\[
\frac{g'(a)}{f'(a)}=\frac{1-|b|^2}{2\sqrt{|b|}(1-|b|)}
 =\frac{1+|b|}{2\sqrt{|b|}}>1,
\]
因为 $0<|b|<1$. 这与最大导数矛盾, 所以像覆盖整个圆盘. 单射全纯逆由局部映射定理全纯.

若 $f_1,f_2$ 均归一化, 则 $f_2\circ f_1^{-1}$ 为固定0的圆盘自同构, 只能是旋转. 导数为正实数迫使旋转为1, 得唯一性.
\end{proof}
不能把 $\Omega=\C$ 纳入, 否则映到圆盘会是有界非常数整函数, 与 Liouville 矛盾. 定理也不保证映射有初等闭式表达, 不自动保证任意粗糙边界都有连续一一的边界延伸. 显式公式与一般存在性解决的是不同问题.

\originalsection{Schwarz 反射}
\begin{theorem}[直线反射原理]
设域关于实轴局部对称, $f$ 在上半部分全纯, 连续到一段实轴且该段函数值为实数. 则可向下半部分延拓为 $F(z)=\overline{f(\bar z)}$, 并在实轴段附近全纯.
\end{theorem}
\begin{proof}
下半部分的共轭差商给 $F'(z)=\overline{f'(\bar z)}$, 所以全纯. 实轴边值为实数使两部分与原边值连续衔接. 对跨实轴的小三角形, 在距轴 $\eps$ 的上下两条平行线截开, 每个远离实轴的多边形边界积分为零. 让 $\eps\to0$, 连续性使沿轴两份积分反向抵消, 连接短边积分趋零, 故原三角形积分为零. Morera 给穿过实轴的全纯性.
\end{proof}
实轴段上没有连续实值边界这一条件时不能直接反射. 圆弧边界可先用 Möbius 变换送到实轴, 对像函数反射后再拉回; 所有映射须在所选局部正常点有效. 反射可帮助延拓, 但不是任意实轴边界函数都能全纯延伸.

\originalsection{习题与解答}
\begin{exercise}
全纯 $f:\D\to\D$ 固定两个不同的内部点. 证明 $f$ 恒等.
\end{exercise}
\begin{solution}
用圆盘自同构把一个固定点送到0, 得共轭映射固定0及另一非零点. Schwarz 在非零固定点处取等, 映射为旋转, 第二固定点又迫使旋转为1. 反向共轭得到原映射恒等.
\end{solution}
\begin{exercise}
证明若 $f_n:\Omega\to\D$ 全纯, 则存在局部一致收敛子序列. 极限一定仍映入开圆盘吗?
\end{exercise}
\begin{solution}
局部统一有界, Montel 给子序列. 极限模至多1; 若非常数, 最大模原理使其严格小于1. 常数极限可以在圆周上, 例如 $f_n=1-1/n$ ($n\ge2$).
\end{solution}
\begin{exercise}
说明穿孔圆盘不能与单位圆盘双全纯.
\end{exercise}
\begin{solution}
若有双全纯 $T$, 它是拓扑同胚, 将每条闭路收缩与其逆一起拉回, 使穿孔圆盘单连通. 但小圆围0绕数为1, 若可收缩成点, 同伦不变性迫使其绕数为0, 矛盾.
\end{solution}

\originalchapter{来源与覆盖}
\originalsection{基础主课}
原件页数包括各文件的官方署名尾页. 下表按实际独立文件计数, 不按课程日历重复链接计数. 主课13专题均编入; 同类机械计算例合并, 不宣称逐句或逐例复刻. 原件无渐显重复页组, 不做原文件删页. 中文删去的是重复介绍、复习句与每份文件相同的署名尾页.
\begin{longtable}{p{.19\textwidth}p{.06\textwidth}p{.64\textwidth}}
\toprule 文件 & 页数 & 中文位置与内容\\\midrule\endhead
\href{https://ocw.mit.edu/courses/18-04-complex-variables-with-applications-spring-2018/resources/mit18_04s18_topic0/}{topic0} & 3 & 前言\\
\href{https://ocw.mit.edu/courses/18-04-complex-variables-with-applications-spring-2018/resources/mit18_04s18_topic1/}{topic1} & 25 & 第1章\\
\href{https://ocw.mit.edu/courses/18-04-complex-variables-with-applications-spring-2018/resources/mit18_04s18_topic2/}{topic2} & 21 & 第2章及第1章球面\\
\href{https://ocw.mit.edu/courses/18-04-complex-variables-with-applications-spring-2018/resources/mit18_04s18_1802review/}{1802review} & 10 & 多元微积分先修说明及第3、9、10章\\
\href{https://ocw.mit.edu/courses/18-04-complex-variables-with-applications-spring-2018/resources/mit18_04s18_topic3/}{topic3} & 11 & 第3章、第8章绕数\\
\href{https://ocw.mit.edu/courses/18-04-complex-variables-with-applications-spring-2018/resources/mit18_04s18_topic4/}{topic4} & 17 & 第3章ML、第4章、第7章实积分\\
\href{https://ocw.mit.edu/courses/18-04-complex-variables-with-applications-spring-2018/resources/mit18_04s18_topic5/}{topic5} & 9 & 第9章\\
\href{https://ocw.mit.edu/courses/18-04-complex-variables-with-applications-spring-2018/resources/mit18_04s18_topic6/}{topic6} & 13 & 第10章\\
\href{https://ocw.mit.edu/courses/18-04-complex-variables-with-applications-spring-2018/resources/mit18_04s18_topic7/}{topic7} & 19 & 第5、6章\\
\href{https://ocw.mit.edu/courses/18-04-complex-variables-with-applications-spring-2018/resources/mit18_04s18_topic8/}{topic8} & 18 & 第6章\\
\href{https://ocw.mit.edu/courses/18-04-complex-variables-with-applications-spring-2018/resources/mit18_04s18_topic9/}{topic9} & 24 & 第7章、第12章Fourier/Green\\
\href{https://ocw.mit.edu/courses/18-04-complex-variables-with-applications-spring-2018/resources/mit18_04s18_topic10/}{topic10} & 22 & 第9--11章\\
\href{https://ocw.mit.edu/courses/18-04-complex-variables-with-applications-spring-2018/resources/mit18_04s18_topic11/}{topic11} & 14 & 第8章\\
\href{https://ocw.mit.edu/courses/18-04-complex-variables-with-applications-spring-2018/resources/mit18_04s18_topic12/}{topic12} & 16 & 第12章\\
\href{https://ocw.mit.edu/courses/18-04-complex-variables-with-applications-spring-2018/resources/mit18_04s18_topic13/}{topic13} & 7 & 第13章\\
\bottomrule\end{longtable}
\originalsection{拓展补充}
18.112的22份实际补充讲义全部登记与归档, 但不是22章教材. 下表明确本册采用、背景及未编内容. 公开补充多为教材页码注记, 不把教材中未提供的证明冒称为公开讲义内容.
\begin{longtable}{p{.19\textwidth}p{.06\textwidth}p{.64\textwidth}}
\toprule 文件 & 页数 & 编入范围\\\midrule\endhead
\href{https://ocw.mit.edu/courses/18-112-functions-of-a-complex-variable-fall-2008/resources/lecture1/}{lecture1} & 4 & 已归档但未编入独立拓展专题\\
\href{https://ocw.mit.edu/courses/18-112-functions-of-a-complex-variable-fall-2008/resources/lecture2/}{lecture2} & 7 & 已归档但未编入独立拓展专题\\
\href{https://ocw.mit.edu/courses/18-112-functions-of-a-complex-variable-fall-2008/resources/lecture3/}{lecture3} & 4 & 已归档但未编入独立拓展专题\\
\href{https://ocw.mit.edu/courses/18-112-functions-of-a-complex-variable-fall-2008/resources/lecture4/}{lecture4} & 3 & 第5章主题与收敛背景; 主源附录未完成, 证明编者补足\\
\href{https://ocw.mit.edu/courses/18-112-functions-of-a-complex-variable-fall-2008/resources/lecture5/}{lecture5} & 2 & 已归档但未编入独立拓展专题\\
\href{https://ocw.mit.edu/courses/18-112-functions-of-a-complex-variable-fall-2008/resources/lecture6/}{lecture6} & 3 & 第11章主题背景; 补充PDF不替代教材全部正文\\
\href{https://ocw.mit.edu/courses/18-112-functions-of-a-complex-variable-fall-2008/resources/lecture7/}{lecture7} & 5 & 第11章广义圆分类与背景; 交比证明编者补足\\
\href{https://ocw.mit.edu/courses/18-112-functions-of-a-complex-variable-fall-2008/resources/lecture8/}{lecture8} & 4 & 已归档但未编入独立拓展专题\\
\href{https://ocw.mit.edu/courses/18-112-functions-of-a-complex-variable-fall-2008/resources/lecture9/}{lecture9} & 2 & 已归档但未编入独立拓展专题\\
\href{https://ocw.mit.edu/courses/18-112-functions-of-a-complex-variable-fall-2008/resources/lecture10/}{lecture10} & 4 & 第4--5章主题背景\\
\href{https://ocw.mit.edu/courses/18-112-functions-of-a-complex-variable-fall-2008/resources/lecture11/}{lecture11} & 3 & 第6章主题背景\\
\href{https://ocw.mit.edu/courses/18-112-functions-of-a-complex-variable-fall-2008/resources/lecture12/}{lecture12} & 5 & 第8、14章局部映射、Schwarz与长度元\\
\href{https://ocw.mit.edu/courses/18-112-functions-of-a-complex-variable-fall-2008/resources/lecture13/}{lecture13} & 5 & 第8章Dixon一般Cauchy证明\\
\href{https://ocw.mit.edu/courses/18-112-functions-of-a-complex-variable-fall-2008/resources/lecture14/}{lecture14} & 6 & 第6、8章留数与辐角主题背景\\
\href{https://ocw.mit.edu/courses/18-112-functions-of-a-complex-variable-fall-2008/resources/lecture15/}{lecture15} & 5 & 第7章轮廓积分背景\\
\href{https://ocw.mit.edu/courses/18-112-functions-of-a-complex-variable-fall-2008/resources/lecture16/}{lecture16} & 6 & 第9章共轭、环域均值、Poisson\\
\href{https://ocw.mit.edu/courses/18-112-functions-of-a-complex-variable-fall-2008/resources/lecture17/}{lecture17} & 5 & 已归档但未编入独立拓展专题\\
\href{https://ocw.mit.edu/courses/18-112-functions-of-a-complex-variable-fall-2008/resources/lecture18_long2/}{lecture18\_long2} & 6 & 第13章Gamma积分、Beta、反射与乘积路线\\
\href{https://ocw.mit.edu/courses/18-112-functions-of-a-complex-variable-fall-2008/resources/lecture19/}{lecture19} & 3 & 第14章Montel紧性与稠密点修正\\
\href{https://ocw.mit.edu/courses/18-112-functions-of-a-complex-variable-fall-2008/resources/lecture20/}{lecture20} & 1 & 第14章Riemann主题; 原件仅1页教材证明注记, 完整证明为编者补充\\
\href{https://ocw.mit.edu/courses/18-112-functions-of-a-complex-variable-fall-2008/resources/lecture21_22/}{lecture21\_22} & 11 & 已归档但未编入独立拓展专题\\
\href{https://ocw.mit.edu/courses/18-112-functions-of-a-complex-variable-fall-2008/resources/lecture23/}{lecture23} & 2 & 已归档但未编入独立拓展专题\\
\bottomrule\end{longtable}
\originalsection{边界与改动}
主课的定义、定理、证明和主要应用按依赖重排. 原课无证明的大 Picard 定理保留为明示引用版, 后文不依赖; Gamma 的复扇区 Stirling 渐近不在本书证明范围, 正实轴版本有完整证明. 反三角函数采用局部分支与复合割线实例, 不枚举每个逆函数的所有分支图. 两分支点积分采用全纯分支与原函数完成计算; 振子 Fourier 主值例明确区别对称 Green 核与因果解, 因果题解完整给出. 这些方法替换均未冒称为逐句原译.
18.112未逐专题译编的内容包括 L1的补充习题、L2部分留给学生的题、L3的 Gauss--Lucas 与特殊有理函数题, L5教材勘误、L8换元附题、L9绕数短注、L17一般 Mittag--Leffler 构造及 zeta级数题, L18的整函数插值拓展, L21--L22素数定理及 L23 zeta延拓/函数方程. 本书的Gamma乘积不等于一般 Weierstrass 规范乘积理论. 完整双曲几何、边界延伸定理与一般亚纯正规族同样不在本册范围.
课程独立作业、习题课、考试、MATLAB教程及外部指定教材没有整批吞入. 本书习题为编者自设并附解答, 不标作 OCW 官方答案. 原讲义内纯文字参考教材的署名保留在原件, 本书未复制教材整页或受限第三方图. 所有正文图为原生 TikZ.
逐文件官方资源页、实际PDF地址、字节数、SHA256、页数及许可记录见 source-manifest.json. 原机构 MIT 与原作者署名在前言和参考文献保留. 原件及本中文改编按 CC BY-NC-SA 4.0, 单独权利声明优先. 版本核查日期2026-10-08.

\begin{thebibliography}{9}
\bibitem{orloff} Jeremy Orloff. \textit{18.04 Complex Variables with Applications, Spring 2018}. MIT OpenCourseWare. \href{https://ocw.mit.edu/courses/18-04-complex-variables-with-applications-spring-2018/pages/lecture-notes/}{官方讲义目录}. 13个专题及2份介绍、复习文件.
\bibitem{helgason} Sigurdur Helgason; lecture notes prepared by Zuoqin Wang. \textit{18.112 Functions of a Complex Variable, Fall 2008}. MIT OpenCourseWare. \href{https://ocw.mit.edu/courses/18-112-functions-of-a-complex-variable-fall-2008/pages/lecture-notes/}{官方补充讲义目录}. 22份文件; 本书采用范围见来源表.
\bibitem{terms} MIT OpenCourseWare. \href{https://ocw.mit.edu/pages/privacy-and-terms-of-use/}{Privacy and Terms of Use}; \href{https://creativecommons.org/licenses/by-nc-sa/4.0/}{CC BY-NC-SA 4.0 International}.
\end{thebibliography}
\end{document}
